% Wersja źródłowa: audyt II, jednostka 15B
\setcounter{section}{15}
\noindent\textbf{Wymagane wyniki.} Wykorzystujemy podstawy skalarnej całki Lebesgue'a, twierdzenia o zbieżności zmajoryzowanej i Fubiniego, twierdzenia Hahna--Banacha i o wykresie domkniętym oraz definicje przestrzeni Banacha i jej przestrzeni sprzężonej.\medskip
% \hyperref[sec:integral-vector15b]{Moduł 15B}; umieszczony po 15A, bez przesuwania \hyperref[sec:wariacyjny16]{rozdziałów 16}--33.
\begingroup
\renewcommand{\thesection}{15B}
\renewcommand{\theHsection}{15B}
\section{Całki Bochnera, Pettisa, Dunforda i Gelfanda oraz miary wektorowe}
\label{sec:integral-vector15b}

W \hyperref[sec:integral-lcs15a]{module 15A} całkowaliśmy \emph{funkcje ciągłe na odcinku} w topologii
przestrzeni lokalnie wypukłej. Teraz dopuszczamy funkcje mierzalne;
różne pojęcia całki wynikają z tego, czy granica istnieje w normie,
czy najpierw całkujemy wszystkie obrazy skalarne, i w jakiej przestrzeni
musi leżeć wektor reprezentujący te całki. Nie powtarzamy konstrukcji
skalarnej całki Lebesgue'a.

W całym module $(S,\Sigma,\mu)$ jest \emph{zupełną przestrzenią z miarą
skończoną}, $\K\in\{\R,\C\}$, zaś $X$ jest przestrzenią Banacha nad
$\K$, chyba że jawnie przyjmiemy inne założenia. Przez $X'$ i $X''$
oznaczamy odpowiednio przestrzeń sprzężoną i podwójnie sprzężoną; $J_X:X\to X''$
jest zanurzeniem kanonicznym danym przez $J_X(x)(x')=x'(x)$.
W przestrzeniach $L^1$ utożsamiamy funkcje równe prawie wszędzie.
W wypowiedziach o wartościach punktowych i całce Riemanna
rozpatrujemy konkretne funkcje, a nie klasy równoważności. Dla $A\in\Sigma$ zapis
$\int_A f(s)\,\mathrm d\mu(s)$ oznacza $\int_S\chi_A(s)f(s)\,\mathrm d\mu(s)$.

\subsection{Mierzalność silna i całka Bochnera}

\begin{definicja}[Mierzalność silna i słaba]
\label{def:measurability15b}
Niech $(S,\Sigma,\mu)$ będzie zupełną przestrzenią z miarą
skończoną, a $X$ przestrzenią Banacha nad $\K\in\{\R,\C\}$.
Funkcję $f:S\to X$ nazywamy \emph{silnie mierzalną}, jeżeli istnieją
funkcje proste $s_n=\sum_{j=1}^{m_n}\chi_{A_{n,j}}x_{n,j}$,
$m_n\ge1$, $A_{n,j}\in\Sigma$, $x_{n,j}\in X$, takie że $s_n(s)\to f(s)$
w normie dla prawie każdego $s\in S$. Nazywamy ją
\emph{słabo mierzalną}, jeżeli $x'\circ f:S\to\K$ jest
mierzalna dla każdego $x'\in X'$. Funkcję $f$ nazywamy
\emph{istotnie ośrodkowo-wartościową}, jeśli poza pewnym zbiorem
miary zero jej wartości leżą w ośrodkowej podprzestrzeni domkniętej
przestrzeni $X$.
\end{definicja}

\begin{twierdzenie}[Twierdzenie Pettisa o mierzalności]
\label{thm:pettis-meas15b}
Niech $(S,\Sigma,\mu)$ będzie zupełną przestrzenią z miarą
skończoną, a $X$ przestrzenią Banacha nad $\K\in\{\R,\C\}$.
Funkcja $f:S\to X$ jest silnie mierzalna wtedy i tylko wtedy, gdy
jest słabo mierzalna oraz istotnie ośrodkowo-wartościowa.
\end{twierdzenie}
\begin{proof}
Jeżeli $s_n\to f$ prawie wszędzie, to dla każdego $x'\in X'$
mamy $x'(s_n)\to x'(f)$, co daje mierzalność słabą.
Wartości wszystkich $s_n$ tworzą zbiór przeliczalny;
jego domknięta otoczka liniowa jest ośrodkową podprzestrzenią,
w której leży $f(s)$ poza wspólnym zbiorem miary zero.

Odwrotnie, po usunięciu zbioru miary zero można przyjąć, że $f$
przyjmuje wartości w ośrodkowej przestrzeni Banacha $Y\subset X$.
Jeżeli $Y=\{0\}$, funkcja jest równa zeru prawie wszędzie,
więc jest silnie mierzalna. Dalej zakładamy $Y\ne0$.
Wybierzmy gęsty ciąg $(y_j)$ w $Y$. Istnieje przeliczalna
rodzina $(\ell_k)$ elementów $X'$ o normie co najwyżej $1$
taka, że dla każdego $y\in Y$ zachodzi
$\|y\|=\sup_k|\ell_k(y)|$: otrzymujemy ją, dobierając
funkcjonały normujące gęsty zbiór sfery jednostkowej $Y$
i przedłużając je za pomocą twierdzenia Hahna--Banacha.
Dokładniej, dla gęstego ciągu $(u_k)$ w sferze jednostkowej
$Y$ wybieramy $\ell_k\in X'$ o normie $1$ i
$\ell_k(u_k)=1$. Dla dowolnego $y\ne0$ oraz $\varepsilon>0$
istnieje $k$ takie, że $\|u_k-y/\|y\|\|<\varepsilon$,
więc $|\ell_k(y)|\ge(1-\varepsilon)\|y\|$.
Przeciwne oszacowanie wynika z $\|\ell_k\|=1$;
dla $y=0$ równość jest natychmiastowa.
W konsekwencji $s\mapsto\|f(s)-y_j\|$ jest mierzalna,
ponieważ jest supremum przeliczalnej rodziny mierzalnych
funkcji skalarnych.

Dla $n\ge1$ definiujemy rozłączne, mierzalne zbiory $B_{n,j}$,
przypisując $s$ pierwszemu indeksowi $j$, dla którego
$\|f(s)-y_j\|<1/n$. Zbiory te pokrywają $S$ z dokładnością
do zbioru miary zero, a funkcja przyjmująca wartość $y_j$
na $B_{n,j}$ przybliża $f$ z błędem mniejszym od $1/n$.
Ponieważ $\mu(S)<\infty$, można wybrać $m_n$ tak, aby
$\mu(\bigcup_{j>m_n}B_{n,j})<2^{-n}$, i przyjąć wartość
$0$ na tym ogonie. Powstaje funkcja prosta $s_n$;
z lematu Borela--Cantellego otrzymujemy $s_n(s)\to f(s)$
prawie wszędzie. Jest to mierzalność silna.
\end{proof}

\begin{definicja}[Całka Bochnera]
\label{def:bochner15b}
Niech $(S,\Sigma,\mu)$ będzie zupełną przestrzenią z miarą
skończoną, a $X$ przestrzenią Banacha nad $\K\in\{\R,\C\}$.
Funkcję silnie mierzalną $f:S\to X$ nazywamy
\emph{całkowalną w sensie Bochnera}, jeśli istnieje ciąg
funkcji prostych $(s_n)$ taki, że
$\int_S\|f(s)-s_n(s)\|\,\mathrm d\mu(s)\to0$.
Wtedy definiujemy
$\int_S f(s)\,\mathrm d\mu(s)
=\lim_{n\to\infty}\int_S s_n(s)\,\mathrm d\mu(s)$,
gdzie całka funkcji prostej $\sum_{j=1}^m\chi_{A_j}x_j$
jest równa $\sum_{j=1}^m\mu(A_j)x_j$ dla parami rozłącznych
$A_j\in\Sigma$ i $x_j\in X$. Każda funkcja prosta ma
takie przedstawienie; przez wspólne rozdrobnienie dwóch
skończonych podziałów (z uwzględnieniem zbioru, na którym
funkcja ma wartość zero) i skończoną addytywność $\mu$
otrzymujemy niezależność tej sumy od przedstawienia.
Przez $L^1(S,\mu;X)$, również $L^1(S;X)$ przy ustalonej
mierze $\mu$, oznaczamy przestrzeń klas funkcji całkowalnych
w sensie Bochnera z normą
$\|f\|_{L^1}=\int_S\|f(s)\|\,\mathrm d\mu(s)$.
Istnienie całki i niezależność od wyboru aproksymacji
uzasadniamy w następnym twierdzeniu.
\end{definicja}

\begin{twierdzenie}[Kryterium całkowalności i oszacowanie]
\label{thm:bochner-criterion15b}
Niech $(S,\Sigma,\mu)$ będzie zupełną przestrzenią z miarą
skończoną, a $X$ przestrzenią Banacha nad $\K\in\{\R,\C\}$.
Funkcja $f:S\to X$ jest całkowalna w sensie Bochnera wtedy
i tylko wtedy, gdy jest silnie mierzalna i
$\int_S\|f(s)\|\,\mathrm d\mu(s)<\infty$.
Całka nie zależy od ciągu przybliżeń, jest liniowa i spełnia
\[
 \left\|\int_A f(s)\,\mathrm d\mu(s)\right\|
 \le\int_A\|f(s)\|\,\mathrm d\mu(s)\qquad(A\in\Sigma).
\]
Jeżeli $f$ jest całkowalna w sensie Bochnera, to dla każdej
przestrzeni Banacha $Y$ nad $\K$ i operatora
$T\in\mathcal L(X,Y)$ funkcja $Tf$ jest całkowalna w sensie
Bochnera oraz dla każdego $A\in\Sigma$ zachodzi
$T(\int_A f(s)\,\mathrm d\mu(s))=\int_A T(f(s))\,\mathrm d\mu(s)$.
\end{twierdzenie}
\begin{proof}
Dla funkcji prostych $h$ mamy
$\|\int_A h(s)\,\mathrm d\mu(s)\|\le\int_A\|h(s)\|\,\mathrm d\mu(s)$.
Jeżeli $\int_S\|f-s_n\|\,\mathrm d\mu\to0$, to
$\int_S\|s_n-s_m\|\,\mathrm d\mu\le
\int_S\|s_n-f\|\,\mathrm d\mu+\int_S\|f-s_m\|\,\mathrm d\mu$.
Stąd całki $s_n$ po $A$ są ciągiem Cauchy'ego w zupełnej
przestrzeni $X$. Takie samo oszacowanie dla dwóch różnych
ciągów aproksymujących $f$ daje niezależność granicy od
aproksymacji. Ograniczenie funkcji prostych do $A$ przez
mnożenie przez $\chi_A$ zachowuje zbieżność błędu w całce,
więc uzasadnia również istnienie całki po $A$.
Jeżeli $s_n$ i $t_n$ aproksymują odpowiednio $f$ i $h$,
to $\alpha s_n+\beta t_n$ aproksymują $\alpha f+\beta h$
dla $\alpha,\beta\in\K$; liniowość wynika z liniowości
całki funkcji prostych, wynikającej ze wspólnego
rozdrobnienia ich podziałów. Ponadto
$|\|f\|-\|s_n\||\le\|f-s_n\|$, stąd całkowalność normy
i zbieżność całek norm, co uzasadnia przejście do granicy
w oszacowaniu całki po $A$.

Jeśli z kolei $f$ jest silnie mierzalna i ma całkowalną normę,
wybierzmy funkcje proste $u_n\to f$ prawie wszędzie.
Funkcja $s_n$ jest równa $u_n$ tam, gdzie
$\|u_n(s)\|\le2\|f(s)\|$, a zeru w pozostałych punktach.
Oba zbiory są mierzalne, zatem $s_n$ jest funkcją prostą,
$\|s_n\|\le2\|f\|$, a $s_n\to f$ prawie wszędzie:
dla $f(s)\ne0$ nierówność definiująca $s_n$ jest
spełniona od pewnego miejsca, a dla $f(s)=0$ mamy
$s_n(s)=0$. Stąd $\|s_n-f\|\le3\|f\|$ i skalarne
twierdzenie o zbieżności zmajoryzowanej zapewnia
$\int_S\|f(s)-s_n(s)\|\,\mathrm d\mu(s)\to0$.
Dla operatora $T$ każda funkcja $Ts_n$ jest prosta,
a silna mierzalność $Tf$ wynika z ciągłości $T$ i silnej
mierzalności $f$. Ponadto
$\int\|Tf-Ts_n\|\le\|T\|\int\|f-s_n\|$, a równość
$T\int s_n=\int Ts_n$ przechodzi do granicy.
\end{proof}

\begin{twierdzenie}[Zbieżność zmajoryzowana Bochnera]
\label{thm:bochner-dct15b}
Niech $(S,\Sigma,\mu)$ będzie zupełną przestrzenią z miarą
skończoną, a $X$ przestrzenią Banacha nad $\K\in\{\R,\C\}$.
Niech $f:S\to X$ będzie funkcją, a $(f_n)_{n\ge1}$ ciągiem
silnie mierzalnych funkcji $f_n:S\to X$ takim, że
$f_n(s)\to f(s)$ w normie dla prawie każdego $s\in S$.
Jeżeli istnieje $g\in L^1(S,\mu)$, $g\ge0$, taka że
$\|f_n(s)\|\le g(s)$ prawie wszędzie dla wszystkich $n$,
to $f$ i $f_n$ są całkowalne w sensie Bochnera oraz
$\int_S\|f_n(s)-f(s)\|\,\mathrm d\mu(s)\to0$, a więc
$\int_S f_n(s)\,\mathrm d\mu(s)\to\int_S f(s)\,\mathrm d\mu(s)$ w normie.
\end{twierdzenie}
\begin{proof}
Granica prawie wszędzie funkcji silnie mierzalnych jest
silnie mierzalna: dla każdego $n$ wybieramy funkcję prostą $v_n$
taką, że $\mu(\{s:\|v_n(s)-f_n(s)\|>1/n\})<2^{-n}$.
Jest to możliwe dzięki zbieżności prawie wszędzie aproksymacji
prostych i skończoności miary. Lemat Borela--Cantellego daje
$\|v_n(s)-f_n(s)\|\to0$ prawie wszędzie, więc $v_n(s)\to f(s)$
prawie wszędzie. Ponadto $\|f\|\le g$ i
$\|f_n-f\|\le2g$ prawie wszędzie. Skalarne twierdzenie
o zbieżności zmajoryzowanej i kryterium całkowalności
Bochnera kończą dowód.
\end{proof}

\begin{twierdzenie}[Twierdzenie Fubiniego dla całki Bochnera]
\label{thm:bochner-fubini15b}
Niech $(S,\Sigma,\mu)$ i $(T,\mathcal T,\nu)$ będą
zupełnymi przestrzeniami z miarami skończonymi, a $X$
przestrzenią Banacha nad $\K\in\{\R,\C\}$.
Niech $f:S\times T\to X$ będzie silnie mierzalna względem
uzupełnienia miary produktowej i spełnia
$\int_{S\times T}\|f(s,t)\|\,\mathrm d(\mu\otimes\nu)(s,t)<\infty$.
Całki na $S\times T$ rozumiemy względem uzupełnienia
miary $\mu\otimes\nu$; ich wartości nie zależą od wyboru
przedstawiciela funkcji z dokładnością do zbioru miary zero.
Wtedy dla $\mu$-prawie każdego $s$ funkcja $f(s,\cdot)$
oraz dla $\nu$-prawie każdego $t$ funkcja $f(\cdot,t)$
są całkowalne w sensie Bochnera. Po przyjęciu wartości $0$
dla całek po przekrojach wyjątkowych obie funkcje otrzymane
przez całkowanie po jednej zmiennej są całkowalne w sensie
Bochnera i
\[
 \int_{S\times T}f\,\mathrm d(\mu\otimes\nu)
 =\int_S\left(\int_Tf(s,t)\,\mathrm d\nu(t)\right)\mathrm d\mu(s)
 =\int_T\left(\int_Sf(s,t)\,\mathrm d\mu(s)\right)\mathrm d\nu(t).
\]
\end{twierdzenie}
\begin{proof}
W pierwszej kolejności redukujemy zagadnienie do funkcji mierzalnej
względem nieuzupełnionej produktowej $\sigma$-algebry. Istotnie,
funkcja $f$ jest prawie wszędzie granicą funkcji prostych mierzalnych
względem uzupełnienia. Każdy z przeliczalnie wielu zbiorów
występujących w tych funkcjach różni się o zbiór zawarty w mierzalnym
zbiorze miary produktowej zero od pewnego zbioru produktowo
mierzalnego. Zastępując te zbiory odpowiednimi przedstawicielami,
otrzymujemy ciąg produktowo mierzalnych funkcji prostych $(u_n)$.
Poza jednym produktowo mierzalnym zbiorem $N$ miary zero
ciąg ten zbiega do $f$. Zbiór punktów, w których $(u_n)$
jest ciągiem Cauchy'ego, jest produktowo mierzalny, gdyż
można go zapisać jako
$\bigcap_{k\ge1}\bigcup_{M\ge1}\bigcap_{i,j\ge M}
\{(s,t):\|u_i(s,t)-u_j(s,t)\|<1/k\}$.
Z zupełności $X$ na tym zbiorze istnieje granica; definiujemy
$\widetilde f$ jako tę granicę, a poza tym zbiorem jako $0$.
Mnożąc $u_n$ przez funkcję charakterystyczną zbioru zbieżności,
otrzymujemy produktowo mierzalne funkcje proste zbieżne
punktowo do $\widetilde f$. W szczególności jej norma jest
produktowo mierzalna. Funkcje $f$ i $\widetilde f$ są równe
poza $N$. Na mocy skalarnego twierdzenia Fubiniego przekroje
$N$ mają miarę zero dla prawie wszystkich wartości każdego
z parametrów. Wystarczy więc wykazać tezę dla $\widetilde f$,
którą dalej oznaczamy przez $f$.

Funkcje postaci $\chi_{A\times B}(s,t)x$, gdzie $A\in\Sigma$,
$B\in\mathcal T$ i $x\in X$, spełniają żądane tożsamości.
Ich skończone kombinacje liniowe są gęste w
$L^1(S\times T;X)$: prostokąty generują produktową
$\sigma$-algebrę, a skończoność miar pozwala, za pomocą
argumentu klasy monotonicznej, aproksymować funkcję charakterystyczną
każdego zbioru produktowo mierzalnego w skalarnej normie $L^1$.
Istotnie, klasa zbiorów, których funkcje charakterystyczne należą
do domknięcia liniowej otoczki funkcji charakterystycznych prostokątów,
zawiera algebrę skończonych sum prostokątów (przez rozbicie na
rozłączne prostokąty). Jest zamknięta na rosnące sumy i malejące
przecięcia ciągów zbiorów, ponieważ przy skończonej mierze ich
funkcje charakterystyczne zbiegają w $L^1$. Twierdzenie o klasie
monotonicznej daje zatem całą produktową $\sigma$-algebrę.
Aproksymacja funkcji prostych daje ten sam wynik dla
funkcji o wartościach w $X$.

Wybierzmy zatem skończone kombinacje prostokątne $g_n$
takie, że
$\int_{S\times T}\|f-g_n\|\,\mathrm d(\mu\otimes\nu)<4^{-n}$.
Skalarne twierdzenie Tonellego daje
\[
 \int_S\sum_{n=1}^\infty
 \left(\int_T\|f(s,t)-g_n(s,t)\|\,\mathrm d\nu(t)\right)
 \mathrm d\mu(s)<\infty.
\]
Silna mierzalność przekrojów dla prawie każdego $s$ wynika
z przekrojów produktowo mierzalnych funkcji prostych
aproksymujących $f$ punktowo poza produktowym zbiorem miary
zero: jego przekroje są zerowe dla prawie każdego $s$.
Skalarne twierdzenie Tonellego dla $\|f\|$ daje całkowalność
normy przekroju dla prawie każdego $s$. Powyższe oszacowanie
sumy błędów daje ponadto $g_n(s,\cdot)\to f(s,\cdot)$
w $L^1(T;X)$.
Zgodnie z kryterium Bochnera przekrój jest całkowalny.
Dla takiego $s$ oznaczmy
$F(s)=\int_T f(s,t)\,\mathrm d\nu(t)$ oraz
$F_n(s)=\int_T g_n(s,t)\,\mathrm d\nu(t)$.
Przyjmujemy $F=0$ na wyjątkowym zbiorze miary zero.
Funkcje $F_n$ (określone na całym $S$ przez całkowanie
$g_n$ po drugiej zmiennej) są proste i $F_n(s)\to F(s)$
prawie wszędzie
na mocy oszacowania całki Bochnera po przekroju. Zatem $F$
jest silnie mierzalna i możemy zastosować Tonellego, otrzymując
\[
 \int_S\|F(s)-F_n(s)\|\,\mathrm d\mu(s)
 \le\int_{S\times T}\|f-g_n\|\,\mathrm d(\mu\otimes\nu)
 <4^{-n}.
\]
Z tego oszacowania i całkowalności $F_n$ wynika całkowalność
$F$ w sensie Bochnera.
Tożsamość dla $g_n$, oszacowanie Bochnera i przejście
do granicy dają
$\int_S F\,\mathrm d\mu
 =\int_{S\times T}f\,\mathrm d(\mu\otimes\nu)$.
Zamiana ról $S$ i $T$ prowadzi do drugiej całki iterowanej
oraz kończy dowód. Wszystkie wnioski dla pierwotnej $f$
są prawdziwe z dokładnością do odpowiednich zbiorów
miary zero.
\end{proof}

\begin{stwierdzenie}[Zgodność z całką Riemanna]
\label{prop:riemann-bochner15b}
Niech $a,b\in\R$, $a<b$, a $X$ będzie przestrzenią Banacha
nad $\K\in\{\R,\C\}$. Jeśli $f:[a,b]\to X$ jest ciągła, to jej całka Riemanna
w przestrzeni Banacha z \hyperref[sec:integral-lcs15a]{modułu 15A} jest równa całce Bochnera
względem miary Lebesgue'a na $[a,b]$. Jeżeli $X\ne\{0\}$, całkowalność Bochnera nie implikuje
całkowalności Riemanna, nawet na odcinku $[0,1]$.
\end{stwierdzenie}
\begin{proof}
Ciągłość na zwartym odcinku daje jednostajną ciągłość.
Funkcje proste, stałe na odcinkach drobnego podziału,
przybliżają $f$ jednostajnie; ich całki Bochnera są
odpowiednimi sumami Riemanna, a zbieżność zachodzi w normie.
Kontrprzykładem jest $f=\chi_{\mathbb Q\cap[0,1]}x_0$
dla niezerowego $x_0\in X$: funkcja jest równa zeru prawie
wszędzie, więc ma zerową całkę Bochnera, lecz nie jest
całkowalna w sensie Riemanna: wybierając w każdym
podprzedziale dodatniej długości punkt wymierny, otrzymujemy
sumę Riemanna równą $x_0$, a wybierając wyłącznie punkty
niewymierne, otrzymujemy sumę równą $0$. Oba wybory są
możliwe przy dowolnie małej średnicy podziału, więc nie
istnieje wspólna granica sum. W przypadku $X=\{0\}$
każda funkcja na odcinku jest całkowalna w obu sensach.
\end{proof}

\subsection{Całki wyznaczane przez wartości skalarne}

\begin{definicja}[Całka Dunforda i całka Pettisa]
\label{def:dunford-pettis15b}
Niech $(S,\Sigma,\mu)$ będzie zupełną przestrzenią z miarą
skończoną, a $X$ przestrzenią Banacha nad $\K\in\{\R,\C\}$.
Niech $f:S\to X$ będzie słabo mierzalna i niech
$x'\circ f\in L^1(S,\mu)$ dla każdego $x'\in X'$.
\emph{Całką Dunforda} funkcji $f$ na $A\in\Sigma$ nazywamy
element $D_A\in X''$ określony równością
$D_A(x')=\int_Ax'(f(s))\,\mathrm d\mu(s)$ dla każdego
$x'\in X'$. Funkcję spełniającą powyższe założenia
nazywamy całkowalną w sensie Dunforda; istnienie $D_A$
uzasadnia następne twierdzenie.
Funkcja $f$ jest \emph{całkowalna w sensie Pettisa}, gdy
dla każdego $A\in\Sigma$ istnieje $x_A\in X$ spełniający
$x'(x_A)=\int_Ax'(f(s))\,\mathrm d\mu(s)$ dla każdego
$x'\in X'$. Funkcjonały z $X'$ rozdzielają punkty $X$
na mocy twierdzenia Hahna--Banacha, więc $x_A$ jest
jednoznacznie wyznaczony. Wtedy
$x_A=\int_A^{\mathrm P} f(s)\,\mathrm d\mu(s)$.
Nie wystarcza wymagać reprezentacji tylko dla $A=S$:
wyjaśni to modyfikacja poniższego przykładu z $c_0$.
\end{definicja}

\begin{twierdzenie}[Istnienie całki Dunforda]
\label{thm:dunford-closedgraph15b}
Niech $(S,\Sigma,\mu)$ będzie zupełną przestrzenią
z miarą skończoną, $\K\in\{\R,\C\}$, a $X$ przestrzenią
Banacha nad $\K$. Niech funkcja $f:S\to X$ będzie słabo
mierzalna i niech $x'\circ f\in L^1(S,\mu)$ dla każdego
$x'\in X'$. Wówczas operator $Q_f:X'\to L^1(S,\mu)$,
$Q_f(x')=x'\circ f$, jest liniowy i ograniczony.
Dla każdego $A\in\Sigma$ wzór
$D_A(x')=\int_Ax'(f(s))\,\mathrm d\mu(s)$
określa element $D_A\in X''$, nawet gdy $f$ nie jest
całkowalna w sensie Bochnera, oraz
$\|D_A\|\le\|Q_f\|$.
\end{twierdzenie}
\begin{proof}
Jeżeli $x_n'\to x'$ w normie $X'$ i
$Q_f(x_n')\to h$ w $L^1$, to pewien podciąg
$Q_f(x_{n_k}')$ zbiega do $h$ prawie wszędzie.
Równocześnie dla każdego $s\in S$ mamy
$x_{n_k}'(f(s))\to x'(f(s))$, więc $h=Q_f(x')$ prawie
wszędzie. Wykres operatora $Q_f$ jest zatem domknięty.
Twierdzenie o wykresie domkniętym dla przestrzeni Banacha
daje $\|Q_f\|<\infty$. Stąd
$|D_A(x')|\le\int_A|x'(f(s))|\,\mathrm d\mu(s)
\le\|Q_f\|\|x'\|$, co definiuje element przestrzeni podwójnie sprzężonej.
\end{proof}

\begin{twierdzenie}[Hierarchia trzech całek]
\label{thm:hierarchy-integrals15b}
Niech $(S,\Sigma,\mu)$ będzie zupełną przestrzenią z miarą
skończoną, a $X$ przestrzenią Banacha nad $\K\in\{\R,\C\}$.
Jeżeli $f:S\to X$ jest całkowalna w sensie Bochnera,
to jest całkowalna w sensie Pettisa i Dunforda, a dla
każdego $A\in\Sigma$ zachodzi
$D_A=J_X(\int_A f(s)\,\mathrm d\mu(s))$.
Funkcja $f:S\to X$ słabo mierzalna i skalarnie całkowalna
(tzn. $x'\circ f\in L^1(S,\mu)$ dla każdego $x'\in X'$) jest
całkowalna w sensie Pettisa wtedy i tylko wtedy, gdy
$D_A\in J_X(X)$ dla każdego $A\in\Sigma$.
Całkowalność w sensie Dunforda nie implikuje w ogólności
całkowalności w sensie Pettisa, a całkowalność w sensie
Pettisa nie implikuje całkowalności w sensie Bochnera.
\end{twierdzenie}
\begin{proof}
Dla funkcji prostej równość
$x'(\int_A f(s)\,\mathrm d\mu(s))=\int_Ax'(f(s))\,\mathrm d\mu(s)$
jest bezpośrednia. Przybliżenie w $L^1(S;X)$
i $|x'(v)|\le\|x'\|\|v\|$ pozwalają przejść do granicy
w przypadku Bochnera. Równoważność w drugiej tezie wynika bezpośrednio z definicji:
$D_A=J_X(x_A)$ oznacza, że $D_A(x')=x'(x_A)$ dla każdego
$x'\in X'$, a istnienie takiego $x_A$ dla wszystkich
$A\in\Sigma$ jest warunkiem całkowalności w sensie Pettisa. Zapowiedziane
kontrprzykłady podajemy niżej.
\end{proof}

\begin{przyklad}[Całka Dunforda w $X''$ poza $J_X(X)$]
\label{ex:dunford-not-pettis15b}
Niech $S=\N=\{1,2,\ldots\}$, $\Sigma=\mathcal P(\N)$ i
$\mu(\{n\})=2^{-n}$. W przestrzeni Banacha $X=c_0$
przyjmijmy $f(n)=2^ne_n$, gdzie $(e_n)$ jest standardową
bazą jednostkową $c_0$. Ponieważ $X'=\ell^1$, dla
$a=(a_n)\in\ell^1$ mamy
$\int_S|a(f(n))|\,\mathrm d\mu(n)
=\sum_{n\ge1}|a_n|<\infty$.
Całka Dunforda ma zatem współrzędne
$D_S=(1,1,\ldots)\in\ell^\infty=X''$ i nie leży w
$J_X(c_0)$. Funkcja nie jest całkowalna w sensie Pettisa
ani Bochnera; co więcej,
$\int_S\|f(n)\|\,\mathrm d\mu(n)=\sum_{n\ge1}1=\infty$.

Rozważmy teraz $S_*=\N\times\{-1,1\}$ z $\sigma$-algebrą
wszystkich podzbiorów i miarą
$\mu_*(\{(n,\eta)\})=2^{-n-1}$ dla $\eta\in\{-1,1\}$.
Niech $h(n,\eta)=\eta 2^ne_n\in c_0$.
Dla $a\in\ell^1$ całka modułu skalaryzacji jest równa
$\sum_n|a_n|$, zatem $h$ jest całkowalna w sensie Dunforda.
Na całym $S_*$ wkłady przeciwnych znaków znoszą się i
$D_{S_*}=0\in J_{c_0}(c_0)$ (szereg jest bezwzględnie
zbieżny po każdej skalaryzacji). Natomiast dla
$A=\N\times\{1\}$ otrzymujemy
$D_A=(1/2,1/2,\ldots)\notin J_{c_0}(c_0)$.
Stąd $h$ nie jest całkowalna w sensie Pettisa, mimo że
całka Dunforda po całym $S_*$ jest reprezentowana
przez wektor z $c_0$.
\end{przyklad}

\begin{przyklad}[Całka Pettisa bez całki Bochnera]
\label{ex:pettis-not-bochner15b}
Niech $S=[0,1]$ z miarą Lebesgue'a, zaś
$X=\ell^2([0,1])$. Dla $t\in[0,1]$ definiujemy $f(t)=e_t$.
Każdy $a\in\ell^2([0,1])$ ma co najwyżej przeliczalnie
wiele niezerowych współrzędnych: dla każdego $n\ge1$
zbiór $\{t:|a_t|\ge 1/n\}$ jest skończony, gdyż w przeciwnym
razie sumy kwadratów po skończonych podzbiorach byłyby
nieograniczone. Suma tych zbiorów zawiera wszystkie
niezerowe współrzędne. Na mocy twierdzenia Riesza o
reprezentacji funkcjonałów w przestrzeni Hilberta każda
skalaryzacja ma postać $t\mapsto\langle f(t),a\rangle$
(przy konwencji liniowości względem pierwszej zmiennej).
Jest ona równa zeru prawie wszędzie i jest mierzalna,
ponieważ miara Lebesgue'a jest zupełna. Całki skalarne po każdym $A\in\Sigma$ znikają,
więc $f$ ma całkę Pettisa równą $0$ na każdym $A$.
Nie jest jednak silnie mierzalna: zbiór jej wartości poza
jakimkolwiek zbiorem miary zero jest nieprzeliczalną rodziną
punktów odległych o $\sqrt 2$, a więc nie może leżeć
w ośrodkowej podprzestrzeni $X$. Całka Bochnera nie istnieje,
choć $\|f(t)\|=1$ dla każdego $t$.
\end{przyklad}

\begin{przyklad}[Silna mierzalność i całka Pettisa bez całki Bochnera]
\label{ex:pettis-strong-not-bochner15b}
Niech $S=\N=\{1,2,\ldots\}$, $\Sigma=\mathcal P(\N)$,
$\mu(A)=\sum_{n\in A}2^{-n}$ dla $A\in\Sigma$ oraz
$X=\ell^2(\N;\K)$, gdzie $\K\in\{\R,\C\}$. Miara $\mu$ jest
zupełna i $\mu(S)=1$. Przestrzeń $X$ jest ośrodkową
przestrzenią Hilberta. Oznaczmy przez $(e_n)$ jej standardową
bazę ortonormalną i zdefiniujmy
$f:S\to X$ wzorem $f(n)=\frac{2^n}{n}e_n$.

\emph{Silna mierzalność.} Dla $N\ge1$ funkcja
$s_N=\sum_{n=1}^N\chi_{\{n\}}\frac{2^n}{n}e_n$ jest prosta.
Dla każdego $k\in S$ i każdego $N\ge k$ mamy $s_N(k)=f(k)$,
zatem $s_N(k)\to f(k)$ w normie w każdym punkcie.
W szczególności $f$ jest silnie mierzalna, a na mocy
twierdzenia Pettisa o mierzalności również słabo mierzalna.

\emph{Całkowalność skalarna.} Dla dowolnego $x'\in X'$
twierdzenie Riesza o reprezentacji funkcjonałów daje
wektor $a=(a_n)\in\ell^2(\N;\K)$ taki, że
$x'(v)=\langle v,a\rangle$ dla $v\in X$.
Przyjmujemy iloczyn skalarny liniowy względem pierwszej
zmiennej, więc $x'(e_n)=\overline{a_n}$ (w przypadku
rzeczywistym sprzężenie niczego nie zmienia). Z definicji
miary atomowej i nierówności Cauchy'ego--Schwarza otrzymujemy
\[
 \int_S|x'(f(n))|\,\mathrm d\mu(n)
 =\sum_{n=1}^{\infty}\frac{|a_n|}{n}
 \le \|a\|_{\ell^2}
       \left(\sum_{n=1}^{\infty}\frac1{n^2}\right)^{1/2}
 <\infty.
\]
Zatem $f$ jest całkowalna w sensie Dunforda.

\emph{Reprezentacja całek Pettisa.} Dla każdego $A\in\Sigma$
definiujemy wektor $v_A=(\chi_A(n)/n)_{n\ge1}$.
Ponieważ $\sum_{n\ge1}|\chi_A(n)/n|^2
\le\sum_{n\ge1}1/n^2<\infty$, mamy $v_A\in X$.
Dla dowolnego $x'\in X'$ reprezentowanego przez $a$ jak wyżej
zachodzi
\[
 \int_A x'(f(n))\,\mathrm d\mu(n)
 =\sum_{n\in A}\frac{\overline{a_n}}{n}
 =\langle v_A,a\rangle=x'(v_A).
\]
Szereg skalarny jest bezwzględnie zbieżny na mocy
poprzedniego oszacowania. Ponieważ równość zachodzi dla
każdego $A$ i każdego $x'$, funkcja $f$ jest całkowalna
w sensie Pettisa, a
$\int_A^{\mathrm P} f(n)\,\mathrm d\mu(n)=v_A$.
W szczególności całka po $S$ jest wektorem
$(1,1/2,1/3,\ldots)\in\ell^2$.

\emph{Brak całkowalności Bochnera.} Mamy $\|e_n\|=1$,
więc
\[
 \int_S\|f(n)\|\,\mathrm d\mu(n)
 =\sum_{n=1}^{\infty}2^{-n}\frac{2^n}{n}
 =\sum_{n=1}^{\infty}\frac1n=\infty.
\]
Z kryterium całkowalności Bochnera wynika, że $f$ nie
jest całkowalna w tym sensie, pomimo silnej mierzalności.
W odróżnieniu od \hyperref[ex:pettis-not-bochner15b]{poprzedniego przykładu}
przeszkodą nie jest brak silnej mierzalności, lecz brak
całkowalności normy. Pokazuje to, że nawet w ośrodkowej
przestrzeni Hilberta całkowalność Pettisa jest słabsza
od całkowalności Bochnera.
\end{przyklad}

\begin{uwaga}[Refleksywność i całka Dunforda]
Jeżeli $X$ jest refleksywna, to $J_X(X)=X''$ i każdy
$D_A$ jest reprezentowany przez wektor z $X$. Zatem każda
słabo mierzalna, skalarnie całkowalna funkcja do refleksywnej
przestrzeni Banacha jest całkowalna w sensie Pettisa.
Nie oznacza to całkowalności Bochnera; przykład z
przestrzeni Hilberta $\ell^2([0,1])$ pozostaje ważny.
\end{uwaga}

\begin{twierdzenie}[Zbieżność zmajoryzowana dla całki Pettisa]
\label{thm:pettis-convergence15b}
Niech $(S,\Sigma,\mu)$ będzie zupełną przestrzenią z miarą
skończoną, a $X$ refleksywną przestrzenią Banacha nad
$\K\in\{\R,\C\}$. Niech $f:S\to X$ będzie funkcją, a
$(f_n)_{n\ge1}$ ciągiem słabo mierzalnych funkcji $f_n:S\to X$
takim, że
$f_n(s)\to f(s)$ w normie dla prawie każdego $s$.
Załóżmy, że $\|f_n(s)\|\le g(s)$ dla wszystkich $n$
i prawie wszystkich $s$, gdzie $g\in L^1(S,\mu)$,
$g\ge0$, oraz że funkcje
$s\mapsto\|f_n(s)-f(s)\|$ są mierzalne.
Wtedy $f_n$ i $f$ są całkowalne w sensie Pettisa i dla
każdego $A\in\Sigma$ zachodzi
$\int_A^{\mathrm P}f_n(s)\,\mathrm d\mu(s)
\to\int_A^{\mathrm P}f(s)\,\mathrm d\mu(s)$ w normie $X$.
\end{twierdzenie}
\begin{proof}
Funkcja $f$ jest słabo mierzalna, bo dla każdego $x'\in X'$
jej skalaryzacja jest granicą mierzalnych $x'\circ f_n$.
Dominacja przez $g$ daje całkowalność skalarną $f_n$ i $f$.
Z refleksywności $X$ oraz twierdzenia o istnieniu całki
Dunforda wynika, że wszystkie te funkcje mają całkę Pettisa.
Całka Pettisa jest liniowa na przestrzeni funkcji
całkowalnych w tym sensie: liniowość całek skalarnych
pokazuje, że kombinacja liniowa wektorów reprezentujących
jest wektorem reprezentującym odpowiednią kombinację funkcji,
a jednoznaczność wynika z twierdzenia Hahna--Banacha.
Dla każdego $A\in\Sigma$ i $x'\in X'$ o normie co najwyżej
$1$ zachodzi
$|x'(\int_A^{\mathrm P}(f_n(s)-f(s))\,\mathrm d\mu(s))|
\le\int_A\|f_n(s)-f(s)\|\,\mathrm d\mu(s)$.
Biorąc supremum po takich $x'$ otrzymujemy oszacowanie
normy różnicy całek. Skalarne twierdzenie o zbieżności
zmajoryzowanej stosuje się do mierzalnych funkcji
$\|f_n-f\|\le2g$; jego prawa strona dąży do zera.
\end{proof}

\subsection{Całka Gelfanda i miary wektorowe}

\begin{definicja}[Całka Gelfanda]
\label{def:gelfand15b}
Niech $(S,\Sigma,\mu)$ będzie zupełną przestrzenią z miarą
skończoną, a $Y$ przestrzenią Banacha nad $\K\in\{\R,\C\}$.
Niech $f:S\to Y'$ będzie funkcją \emph{słabo-* mierzalną}, tzn.
$s\mapsto f(s)(y)$ jest mierzalna dla każdego $y\in Y$.
Zakładamy, że $f(\cdot)(y)\in L^1(S,\mu)$ dla każdego
$y\in Y$. \emph{Całką Gelfanda} na $A\in\Sigma$ nazywamy
jedyny $g_A\in Y'$ spełniający
$g_A(y)=\int_Af(s)(y)\,\mathrm d\mu(s)$ dla wszystkich
$y\in Y$. Stosujemy oznaczenie
$g_A=\int_A^{\mathrm G}f(s)\,\mathrm d\mu(s)$.
Istnienie tego funkcjonału uzasadnia następne twierdzenie.
\end{definicja}

\begin{twierdzenie}[Istnienie całki Gelfanda i jej addytywność]
\label{thm:gelfand-exists15b}
Niech $(S,\Sigma,\mu)$ będzie zupełną przestrzenią
z miarą skończoną, $\K\in\{\R,\C\}$, a $Y$
przestrzenią Banacha nad $\K$. Niech $f:S\to Y'$
będzie słabo-* mierzalna oraz
$f(\cdot)(y)\in L^1(S,\mu)$ dla każdego $y\in Y$.
Wtedy dla każdego $A\in\Sigma$ istnieje dokładnie jeden
$g_A\in Y'$ taki, że
$g_A(y)=\int_A f(s)(y)\,\mathrm d\mu(s)$
dla wszystkich $y\in Y$. Odwzorowanie
$A\mapsto g_A$ jest przeliczalnie addytywne
w topologii słabej-* $\sigma(Y',Y)$; bez dodatkowych
założeń nie musi być przeliczalnie addytywne w normie $Y'$.
\end{twierdzenie}
\begin{proof}
Operator $Q:Y\to L^1(S,\mu)$, $(Qy)(s)=f(s)(y)$,
jest domknięty. Istotnie, z $y_n\to y$ i $Qy_n\to h$
w $L^1$ wybieramy podciąg zbieżny prawie wszędzie;
punktowa ciągłość każdego $f(s)\in Y'$ daje $h=Qy$.
Z twierdzenia o wykresie domkniętym $Q$ jest ograniczony.
Zatem $|\int_Af(s)(y)\,\mathrm d\mu(s)|\le\|Q\|\|y\|$,
a wzór definiuje $g_A\in Y'$.
Dla parami rozłącznych $A_n$ i $y\in Y$ skalarna
addytywność całki Lebesgue'a daje
$g_{\bigcup_nA_n}(y)=\sum_ng_{A_n}(y)$.
Jest to dokładnie zbieżność w topologii $\sigma(Y',Y)$.
\end{proof}

\begin{przyklad}[Słaba-* addytywność bez addytywności normowej]
\label{ex:gelfand-weakstar15b}
Niech $Y=\ell^1$, $Y'=\ell^\infty$, a na $S=\N$ przyjmijmy
$\mu(\{n\})=2^{-n}$ oraz $f(n)=2^ne_n\in\ell^\infty$.
Dla $y=(y_n)\in\ell^1$ zachodzi
$\int_S|f(n)(y)|\,\mathrm d\mu(n)=\sum_n|y_n|<\infty$.
W konsekwencji $g_A=\chi_A\in\ell^\infty$, gdzie
$\chi_A=(\chi_A(n))_{n\ge1}$. Dla $A_m=\{m,m+1,\ldots\}$
mamy $\|g_{A_m}\|_\infty=1$ przy każdym $m$, mimo że
$g_{A_m}\to0$ w topologii $\sigma(\ell^\infty,\ell^1)$:
dla $y\in\ell^1$ zachodzi
$|g_{A_m}(y)|\le\sum_{n\ge m}|y_n|\to0$.
Ponieważ $g_S-\sum_{n=1}^N g_{\{n\}}=g_{A_{N+1}}$,
szereg wartości na singletonach nie zbiega w normie do $g_S$.
\end{przyklad}

\begin{definicja}[Miara wektorowa i całkowita wariacja]
\label{def:vector-measure15b}
Niech $(S,\Sigma)$ będzie przestrzenią mierzalną, a $X$
przestrzenią Banacha nad $\K\in\{\R,\C\}$.
Odwzorowanie $m:\Sigma\to X$ nazywamy \emph{miarą wektorową},
jeśli $m(\emptyset)=0$ i dla każdego ciągu parami rozłącznych
$A_n\in\Sigma$ zachodzi
$m(\bigcup_{n\ge1}A_n)=\lim_{N\to\infty}\sum_{n=1}^Nm(A_n)$
w normie $X$. Jej \emph{całkowitą wariacją} nazywamy funkcję
$|m|:\Sigma\to[0,\infty]$ określoną dla $A\in\Sigma$ przez
$|m|(A)=\sup\{\sum_{j=1}^k\|m(A_j)\|:\ A_j\text{ są
parami rozłącznymi zbiorami mierzalnymi, }
\bigcup_{j=1}^kA_j=A,\ k\ge1\}$.
Dopuszczamy puste składniki podziału. W następnym stwierdzeniu
wykażemy, że $|m|$ jest miarą dodatnią.
\end{definicja}

\begin{stwierdzenie}[Przeliczalna addytywność całkowitej wariacji]
\label{prop:variation-measure15b}
Niech $(S,\Sigma)$ będzie przestrzenią mierzalną, a $X$
przestrzenią Banacha nad $\K\in\{\R,\C\}$.
Jeżeli $m:\Sigma\to X$ jest miarą wektorową, to
$|m|:\Sigma\to[0,\infty]$ jest miarą dodatnią (niekoniecznie
skończoną). W szczególności dla parami rozłącznych
$A_n\in\Sigma$ zachodzi
$|m|(\bigcup_{n\ge1}A_n)=\sum_{n\ge1}|m|(A_n)$.
\end{stwierdzenie}
\begin{proof}
Z definicji $|m|(\emptyset)=0$. Niech
$A=\bigcup_{n\ge1}A_n$ będzie sumą parami rozłącznych
zbiorów mierzalnych. Dla każdego $N$ można zestawić
skończone partycje zbiorów $A_1,\ldots,A_N$ w jedną
partycję $A$, dodając ewentualnie
$A\setminus\bigcup_{n=1}^N A_n$.
Po przejściu do supremów i następnie z $N\to\infty$
dostajemy $|m|(A)\ge\sum_{n\ge1}|m|(A_n)$,
również gdy prawa strona jest nieskończona.

Dla nierówności odwrotnej wybierzmy dowolną skończoną
partycję $B_1,\ldots,B_k$ zbioru $A$.
Z przeliczalnej addytywności $m$ w normie i nierówności
trójkąta wynika
\[
 \sum_{j=1}^k\|m(B_j)\|
 \le\sum_{j=1}^k\sum_{n=1}^\infty
        \|m(B_j\cap A_n)\|
 \le\sum_{n=1}^\infty |m|(A_n).
\]
Po wzięciu supremum po wszystkich partycjach $A$
otrzymujemy nierówność odwrotną i tezę.
\end{proof}

\begin{twierdzenie}[Miara generowana przez całkę Bochnera]
\label{thm:bochner-measure15b}
Niech $(S,\Sigma,\mu)$ będzie zupełną przestrzenią z miarą
skończoną, a $X$ przestrzenią Banacha nad $\K\in\{\R,\C\}$.
Jeśli $f:S\to X$ jest całkowalna w sensie Bochnera, to
odwzorowanie $m_f:\Sigma\to X$ określone wzorem
$m_f(A)=\int_Af(s)\,\mathrm d\mu(s)$ jest miarą wektorową,
a jej całkowita wariacja spełnia
\[
 |m_f|(A)=\int_A\|f(s)\|\,\mathrm d\mu(s)
 \qquad(A\in\Sigma).
\]
W szczególności $|m_f|$ jest skończoną miarą bezwzględnie
ciągłą względem $\mu$.
\end{twierdzenie}
\begin{proof}
Dla parami rozłącznych $A_n$ i $A=\bigcup_nA_n$
z oszacowania całki otrzymujemy
$\|m_f(A)-\sum_{n=1}^Nm_f(A_n)\|
\le\int_{\bigcup_{n>N}A_n}\|f(s)\|\,\mathrm d\mu(s)\to0$.
Dla każdej skończonej partycji $A$ suma norm całek
jest ograniczona przez $\int_A\|f\|$, czyli
$|m_f|(A)\le\int_A\|f\|$.
Gdy $f=\sum_j\chi_{B_j}v_j$ jest prosta o rozłącznych
zbiorach poziomicowych, partycja $A$ przez $A\cap B_j$
daje równość. Dla ogólnej $f$ wybierzmy funkcje proste
$f_n\to f$ w $L^1$. Z nierówności
$\bigl||m_f|(A)-|m_{f_n}|(A)\bigr|
\le\int_A\|f(s)-f_n(s)\|\,\mathrm d\mu(s)$
oraz zbieżności norm w $L^1$ otrzymujemy równość.
Istotnie, dla każdej skończonej partycji $A=\bigcup_j A_j$
mamy
$\sum_j\|m_f(A_j)\|\le
\sum_j\|m_{f_n}(A_j)\|+
\sum_j\|m_{f-f_n}(A_j)\|
\le |m_{f_n}|(A)+\int_A\|f-f_n\|\,\mathrm d\mu$.
Biorąc supremum i zamieniając role $f$ oraz $f_n$,
otrzymujemy przytoczone oszacowanie.
\end{proof}

\begin{przyklad}[Miara wektorowa o nieskończonej całkowitej wariacji]
\label{ex:pettis-infinite-variation15b}
Niech $S=\N=\{1,2,\ldots\}$, $\Sigma=\mathcal P(\N)$,
$\mu(\{n\})=2^{-n}$ oraz $X=\ell^2(\N;\K)$,
gdzie $\K\in\{\R,\C\}$. W
\hyperref[ex:pettis-strong-not-bochner15b]{przykładzie dotyczącym silnej mierzalności}
wykazaliśmy, że funkcja $f(n)=\frac{2^n}{n}e_n$ jest
całkowalna w sensie Pettisa, a dla każdego $A\in\Sigma$
jej całka ma postać
\[
 m(A)=\int_A^{\mathrm P}f(n)\,\mathrm d\mu(n)
 =\left(\frac{\chi_A(n)}{n}\right)_{n\ge1}.
\]
Wykażemy, że $m:\Sigma\to X$ jest miarą wektorową
bezwzględnie ciągłą względem $\mu$, lecz jej całkowita
wariacja nie jest skończona.

\emph{Przeliczalna addytywność w normie.}
Oczywiście $m(\emptyset)=0$. Niech $(A_j)_{j\ge1}$ będzie
ciągiem parami rozłącznych zbiorów z $\Sigma$ i niech
$A=\bigcup_{j\ge1}A_j$, $B_N=\bigcup_{j>N}A_j$.
Ponieważ zbiory $A_j$ są rozłączne, dla każdego $n\in\N$
mamy $\chi_A(n)-\sum_{j=1}^N\chi_{A_j}(n)=\chi_{B_N}(n)$.
Zatem
\[
 \left\|m(A)-\sum_{j=1}^Nm(A_j)\right\|_{\ell^2}^2
 =\sum_{n=1}^{\infty}\frac{\chi_{B_N}(n)}{n^2}.
\]
Zbiory $B_N$ maleją do zbioru pustego: każdy element $A$
należy do dokładnie jednego $A_j$, a element spoza $A$
nie należy do żadnego $B_N$. Stąd dla każdego $n$
zachodzi $\chi_{B_N}(n)\to0$. Ponadto
$0\le\chi_{B_N}(n)/n^2\le1/n^2$ i szereg
$\sum_{n\ge1}1/n^2$ jest zbieżny. Twierdzenie o zbieżności
zmajoryzowanej dla miary liczącej daje zbieżność prawej
strony do zera. Otrzymujemy przeliczalną addytywność
$m$ w normie $X$. Jeżeli $\mu(A)=0$, to $A=\emptyset$,
gdyż każdy singleton ma dodatnią miarę; w szczególności
$m(A)=0$.

\emph{Całkowita wariacja.}
Dla dowolnego $A\in\Sigma$ i każdego skończonego podziału
$A=\bigcup_{j=1}^k C_j$ zachodzi
\[
 \sum_{j=1}^k\|m(C_j)\|_{\ell^2}
 =\sum_{j=1}^k\left(\sum_{n\in C_j}\frac1{n^2}\right)^{1/2}
 \le\sum_{j=1}^k\sum_{n\in C_j}\frac1n
 =\sum_{n\in A}\frac1n.
\]
Nierówność wynika z oszacowania normy $\ell^2$ przez
sumę modułów współrzędnych, które dla szeregu nieskończonego
otrzymujemy przez przejście do granicy w oszacowaniu
skończonych sum; dopuszczamy wartość $+\infty$.
Biorąc supremum po podziałach, otrzymujemy
$|m|(A)\le\sum_{n\in A}1/n$.
Z drugiej strony, dla każdego $N\ge1$ możemy podzielić $A$
na singletony $\{n\}$ dla $n\in A\cap\{1,\ldots,N\}$
oraz pozostały zbiór $A\setminus\{1,\ldots,N\}$.
Ponieważ $m(\{n\})=e_n/n$, definicja wariacji daje
$|m|(A)\ge\sum_{n\in A,\ n\le N}1/n$.
Przechodząc z $N\to\infty$, otrzymujemy
\[
 |m|(A)=\sum_{n\in A}\frac1n
 \qquad(A\in\Sigma).
\]
W szczególności $|m|(S)=\sum_{n\ge1}1/n=\infty$.

Zauważmy, że wartości samej miary są jednostajnie ograniczone:
$\|m(A)\|_{\ell^2}\le(\sum_{n\ge1}1/n^2)^{1/2}$
dla każdego $A$. Mimo to jej całkowita wariacja jest
nieskończona. Przykład pokazuje, że miara generowana przez
całkę Pettisa może być przeliczalnie addytywna w normie,
lecz nie spełniać warunku skończonej wariacji, który
zawsze zachodzi dla miary generowanej przez całkę Bochnera.
\end{przyklad}

\subsection{Własności całek wektorowych i zależność od zbioru całkowania}

W dalszym ciągu miara jest skończona i zupełna. Zasadnicza
różnica pomiędzy całkami nie polega na liniowości, która
zachodzi dla wszystkich czterech całek, lecz na przestrzeni
wartości, rodzinie funkcjonałów użytych do skalaryzacji
oraz topologii, w której można sumować całki po rozłącznych
zbiorach. Wprowadzamy oznaczenia $P_A(f)=\int_A^{\mathrm P}f\,\mathrm d\mu$,
$D_A(f)$ dla całki Dunforda i $G_A(h)=\int_A^{\mathrm G}h\,\mathrm d\mu$.

\begin{twierdzenie}[Liniowość, ograniczenie i mnożenie przez funkcje skalarne]
\label{thm:vector-integral-multipliers15b}
Niech $(S,\Sigma,\mu)$ będzie zupełną przestrzenią z miarą
skończoną, $X,Y$ przestrzeniami Banacha nad $\K\in\{\R,\C\}$,
$f:S\to X$ funkcją całkowalną w sensie Dunforda, a
$h:S\to Y'$ funkcją słabo-* mierzalną, dla której
$h(\cdot)(y)\in L^1(S,\mu)$ dla każdego $y\in Y$.
Wprowadźmy dla $A\in\Sigma$
\[
 \begin{split}
 q_f(A)&=\sup_{\|x'\|\le1,\ x'\in X'}
       \int_A|x'(f(s))|\,\mathrm d\mu(s),\\
 r_h(A)&=\sup_{\|y\|\le1,\ y\in Y}
       \int_A|h(s)(y)|\,\mathrm d\mu(s).
 \end{split}
\]
Wówczas $q_f(S)$ i $r_h(S)$ są skończone. Dla każdej
$\varphi\in L^\infty(S,\mu;\K)$ (wybieramy mierzalnego
przedstawiciela klasy) funkcja $\varphi f$ jest
całkowalna w sensie Dunforda, $\varphi h$ w sensie Gelfanda,
a gdy $f$ jest całkowalna w sensie Pettisa, to także
$\varphi f$ jest całkowalna w sensie Pettisa. Dla każdego
$A\in\Sigma$ zachodzą oszacowania
\[
 \begin{split}
 \|D_A(\varphi f)\|&\le\|\varphi\|_\infty q_f(A),\\
 \|G_A(\varphi h)\|&\le\|\varphi\|_\infty r_h(A),\\
 \|P_A(\varphi f)\|&\le\|\varphi\|_\infty q_f(A)
       \quad\text{(w przypadku Pettisa).}
 \end{split}
\]
Każda z tych całek jest liniowa względem funkcji
całkowanej w odpowiedniej klasie całkowalności i nie zmienia
się po zmianie funkcji na zbiorze miary zero. Dla ustalonego
$A$ jest też liniowa względem mnożnika $\varphi$.
Jeżeli $f$ jest całkowalna w sensie Bochnera, to
$\varphi f$ jest całkowalna w tym sensie i
$q_f(A)\le\int_A\|f\|\,\mathrm d\mu$.
\end{twierdzenie}
\begin{proof}
Z \hyperref[thm:dunford-closedgraph15b]{twierdzenia o istnieniu całki Dunforda}
mamy $q_f(S)=\|Q_f\|$, gdzie $Q_f(x')=x'\circ f$.
Analogicznie dowód
\hyperref[thm:gelfand-exists15b]{twierdzenia o istnieniu całki Gelfanda}
daje $r_h(S)=\|Q\|$, gdzie $(Qy)(s)=h(s)(y)$.
Mnożenie funkcji skalarnej z $L^1$ przez funkcję z
$L^\infty$ zachowuje całkowalność, a wszystkie
odpowiednie skalaryzacje pozostają mierzalne. Definicje
całek Dunforda i Gelfanda dają więc istnienie oraz
oszacowania, gdyż dla $\|x'\|\le1$ mamy
$|\int_A\varphi(s)x'(f(s))\,\mathrm d\mu(s)|
\le\|\varphi\|_\infty\int_A|x'(f(s))|\,\mathrm d\mu(s)$,
a dla Gelfanda stosujemy ten sam argument do $h(s)(y)$.
Bierzemy supremum po odpowiedniej kuli jednostkowej.

Niech teraz $f$ będzie całkowalna w sensie Pettisa.
Dla skalarnej funkcji prostej
$\psi=\sum_{j=1}^k c_j\chi_{B_j}$ wektor
$\sum_j c_j P_{A\cap B_j}(f)$ reprezentuje całki skalarne
funkcji $\psi f$ po $A$, co wynika z liniowości
całki skalarnej. Zatem $\psi f$ jest całkowalna w sensie
Pettisa. Z twierdzenia Hahna--Banacha i poprzedniego
oszacowania wynika
$\|P_A(\psi f)\|\le\|\psi\|_\infty q_f(A)$.
Wybierzmy skalarne funkcje proste $\psi_k$ zbieżne do
$\varphi$ w normie $L^\infty$ (po wyborze ograniczonego
przedstawiciela wystarczy podzielić jego ograniczony
zbiór wartości na zbiory o coraz mniejszej średnicy).
Z liniowości dla funkcji prostych dostajemy
\[
 \|P_A(\psi_k f)-P_A(\psi_l f)\|
 \le\|\psi_k-\psi_l\|_\infty q_f(A).
\]
Zupełność $X$ daje granicę $v_A$ w normie. Dla każdego
$x'\in X'$ różnica całek skalarnych odpowiadających
$\psi_k f$ i $\varphi f$ ma moduł nie większy niż
$\|\psi_k-\varphi\|_\infty\int_A|x'(f)|\,\mathrm d\mu$,
a więc dąży do zera. Ciągłość $x'$ daje
$x'(v_A)=\int_A\varphi x'(f)\,\mathrm d\mu$.
Ponieważ $A$ było dowolne, jest to całkowalność Pettisa.
Przejście do granicy daje także oszacowanie normy.

Liniowość i niezmienniczość względem zmian na zbiorach
miary zero wynikają z odpowiednich własności całki
skalarnej oraz jednoznaczności reprezentacji: w $X$ przez
funkcjonały z $X'$, w $X''$ przez wartości na $X'$, a
w $Y'$ przez wartości na $Y$. W przypadku Bochnera
$\varphi f$ jest silnie mierzalna (iloczyny prostych
aproksymacji punktowo przybliżają iloczyn poza zbiorem
zerowym), a $\|\varphi f\|\le\|\varphi\|_\infty\|f\|$
prawie wszędzie. Kryterium Bochnera kończy ten przypadek;
nierówność dla $q_f$ wynika z $|x'(f)|\le\|x'\|\|f\|$.
\end{proof}

\begin{stwierdzenie}[Przechodzenie przez operatory liniowe]
\label{prop:vector-integrals-operators15b}
Niech $(S,\Sigma,\mu)$ będzie zupełną przestrzenią z miarą
skończoną, $X,Z,Y,W$ przestrzeniami Banacha nad
$\K\in\{\R,\C\}$, $T\in\mathcal L(X,Z)$ i
$R\in\mathcal L(W,Y)$. Jeżeli $f:S\to X$ jest
całkowalna w sensie Dunforda, to $Tf$ jest całkowalna
w sensie Dunforda i dla każdego $A\in\Sigma$
$D_A(Tf)=T''D_A(f)$.
Jeżeli $f$ jest całkowalna w sensie Pettisa, to $Tf$
jest całkowalna w sensie Pettisa i
$P_A(Tf)=T P_A(f)$.
Jeżeli $h:S\to Y'$ jest całkowalna w sensie Gelfanda,
to $R'h:S\to W'$ jest całkowalna w sensie Gelfanda i
$G_A(R'h)=R'G_A(h)$.
\end{stwierdzenie}
\begin{proof}
Przez $T':Z'\to X'$ rozumiemy operator
$T'z'=z'\circ T$, a przez $T'':X''\to Z''$ jego
operator sprzężony. Dla $z'\in Z'$ skalaryzacja
$z'\circ Tf=(T'z')\circ f$ jest mierzalna i całkowalna.
Zatem
$D_A(Tf)(z')=D_A(f)(T'z')=(T''D_A(f))(z')$.
Jeżeli istnieje wektor Pettisa $P_A(f)$, to
$z'(TP_A(f))=(T'z')(P_A(f))
=\int_Az'(Tf(s))\,\mathrm d\mu(s)$.
To dowodzi zarówno całkowalności, jak i żądanej równości.
Dla Gelfanda, $w\in W$ i $R':Y'\to W'$ mamy
$(R'h(s))(w)=h(s)(Rw)$. Ta funkcja skalarna jest
mierzalna i całkowalna na mocy założeń, a jej całka
jest równa $G_A(h)(Rw)=(R'G_A(h))(w)$.
Funkcjonały są równe, ponieważ przyjmują te same wartości
na każdym $w\in W$.
\end{proof}

\begin{stwierdzenie}[Addytywność Dunforda i Gelfanda]
\label{prop:dunford-gelfand-additivity15b}
Niech $(S,\Sigma,\mu)$ będzie zupełną przestrzenią z miarą
skończoną, $X,Y$ przestrzeniami Banacha nad
$\K\in\{\R,\C\}$, $f:S\to X$ funkcją całkowalną
w sensie Dunforda i $h:S\to Y'$ funkcją całkowalną
w sensie Gelfanda. Odwzorowania $A\mapsto D_A(f)$
i $A\mapsto G_A(h)$ są przeliczalnie addytywne odpowiednio
w topologiach $\sigma(X'',X')$ i $\sigma(Y',Y)$,
lecz w ogólności nie w normie. Są skończenie addytywne
w normie, jednostajnie ograniczone oraz zerują się na
zbiorach miary zero.
Jeżeli istnieje $g\in L^1(S,\mu)$, $g\ge0$, taka że
$\|f(s)\|\le g(s)$ prawie wszędzie, to
$A\mapsto D_A(f)$ jest miarą wektorową o wariacji
$|D|(A)\le\int_A g\,\mathrm d\mu$.
Analogicznie, jeżeli $\|h(s)\|\le g(s)$ prawie wszędzie,
to $A\mapsto G_A(h)$ jest miarą wektorową
o wariacji $|G|(A)\le\int_A g\,\mathrm d\mu$.
Nie wymaga się mierzalności samych norm.
\end{stwierdzenie}
\begin{proof}
Dla ciągu parami rozłącznych $A_j\in\Sigma$ oraz
$A=\bigcup_j A_j$ skalarna addytywność daje
$D_A(f)(x')=\sum_j D_{A_j}(f)(x')$ dla każdego $x'\in X'$,
a dla Gelfanda $G_A(h)(y)=\sum_jG_{A_j}(h)(y)$ dla każdego
$y\in Y$. Są to dokładnie deklarowane zbieżności słabe-*.
Skończona addytywność jest równością wektorów, gdyż
w obu przypadkach punkty są rozdzielane przez wskazane
skalaryzacje. Oszacowania z
\hyperref[thm:vector-integral-multipliers15b]{twierdzenia o mnożnikach}
dają jednostajne ograniczenie, a zerowanie się na zbiorach
miary zero wynika z własności całek skalarnych.
Brak addytywności w normie Dunforda pokazuje pierwszy
\hyperref[ex:dunford-not-pettis15b]{przykład z $c_0$}:
$D_S=(1,1,\ldots)$, a $\sum_{n=1}^ND_{\{n\}}$ ma jedynki
tylko na pierwszych $N$ pozycjach, zatem norma różnicy
wynosi $1$. Dla Gelfanda analogiczny argument podano
w \hyperref[ex:gelfand-weakstar15b]{przykładzie z $\ell^\infty$}.

Przy dominacji przez $g$ mamy
$\|D_B(f)\|\le\int_Bg\,\mathrm d\mu$ dla każdego $B$,
ponieważ $|x'(f)|\le g$ dla $\|x'\|\le1$ prawie wszędzie.
Stosując skończoną addytywność do $A$ i reszty
$B_N=\bigcup_{j>N}A_j$, dostajemy
$\|D_A(f)-\sum_{j=1}^ND_{A_j}(f)\|\le\int_{B_N}g\,\mathrm d\mu$.
Prawa strona dąży do zera z twierdzenia o zbieżności
zmajoryzowanej, ponieważ $B_N\downarrow\emptyset$.
Dla każdego skończonego podziału $A$ suma norm całek
jest ograniczona przez $\int_Ag\,\mathrm d\mu$;
supremum daje oszacowanie wariacji. Dowód dla Gelfanda
jest identyczny po zastąpieniu $x'(f)$ przez $h(\cdot)(y)$.
\end{proof}

\begin{twierdzenie}[Twierdzenie Orlicza--Pettisa]
\label{thm:orlicz-pettis15b}
Niech $X$ będzie przestrzenią Banacha nad $\K\in\{\R,\C\}$
i niech $(u_n)_{n\ge1}$ będzie ciągiem w $X$.
Załóżmy, że dla każdego $M\subseteq\N$ istnieje $v_M\in X$
taki, że $\sum_{n\in M,\ n\le N}u_n\to v_M$ w topologii
$\sigma(X,X')$ przy $N\to\infty$. Wtedy szereg
$\sum_{n\ge1}u_n$ jest bezwarunkowo zbieżny w normie,
a dla każdego $M$ suma podszeregu w normie jest równa $v_M$.
\end{twierdzenie}
\begin{proof}
Dla każdego $x'\in X'$ każdy podszereg
$\sum_{n\in M}x'(u_n)$ jest zbieżny. W konsekwencji
$\sum_n|x'(u_n)|<\infty$: w przypadku rzeczywistym
w przeciwnym razie podszereg wyrazów nieujemnych
lub niedodatnich byłby rozbieżny. W przypadku zespolonym
stosujemy to do części rzeczywistych i urojonych
(ich podszeregi są zbieżne), a następnie
$|z|\le|\operatorname{Re}z|+|\operatorname{Im}z|$.

Niech $E=\{0,1\}^{\N}$ z topologią produktową, a dla
$\alpha=(\alpha_n)\in E$ niech
$V(\alpha)=v_{\{n:\alpha_n=1\}}$. Przestrzeń $E$ jest
zupełna względem metryki
$d(\alpha,\beta)=\sum_n2^{-n}|\alpha_n-\beta_n|$:
ciąg Cauchy'ego stabilizuje się na każdej współrzędnej,
a zbieżność do ciągu tych ustalonych wartości wynika
z oszacowania ogona przez $\sum_{n>N}2^{-n}$.
Dla każdego $x'\in X'$ mamy
$x'(V(\alpha))=\sum_n\alpha_n x'(u_n)$. Szereg ten
zbiega jednostajnie na $E$, ponieważ jego ogon jest
ograniczony przez $\sum_{n>N}|x'(u_n)|$.
Sumy skończone są ciągłe, więc $x'\circ V$ jest ciągła.
To oznacza ciągłość $V$ z $E$ do $(X,\sigma(X,X'))$.

Niech $Z=\overline{\lin\{u_n:n\ge1\}}$ (domknięcie w normie).
Jest to ośrodkowa podprzestrzeń: kombinacje skończone
o współczynnikach wymiernych (w przypadku zespolonym
o wymiernych częściach rzeczywistych i urojonych)
tworzą w niej zbiór przeliczalny i gęsty. Podprzestrzeń
$Z$ jest słabo domknięta z twierdzenia Hahna--Banacha:
dla punktu spoza $Z$ istnieje ciągły funkcjonał zerujący
się na $Z$, lecz nie w tym punkcie. Wszystkie sumy
skończone podszeregów leżą w $Z$, a więc również ich
słabe granice $V(\alpha)$ leżą w $Z$.

Ustalmy $\varepsilon>0$ i wybierzmy gęsty ciąg $(z_j)$
w $Z$. Zbiory
$F_j=\{\alpha\in E:\|V(\alpha)-z_j\|\le\varepsilon/3\}$
pokrywają $E$. Są domknięte w $E$: domknięta kula
w $X$ jest słabo domknięta, ponieważ normę wyznacza
supremum modułów funkcjonałów z kuli jednostkowej $X'$,
a $V$ jest słabo ciągłe. Twierdzenie Baire'a daje indeks
$j$, dla którego $F_j$ ma niepuste wnętrze.
Bazę topologii produktowej tworzą zbiory określone przez
ustalenie skończonej liczby współrzędnych. Istnieją zatem
$N_0$ i wartości $\alpha_1,\ldots,\alpha_{N_0}$ takie,
że wszystkie ciągi o tych początkowych współrzędnych
należą do $F_j$.

Dla dowolnego skończonego
$H\subseteq\{N_0+1,N_0+2,\ldots\}$ wybierzmy dwa ciągi
z tego zbioru bazowego: pierwszy po $N_0$ ma jedynki
dokładnie na $H$, drugi same zera. Odpowiednie wartości
$V$ różnią się o $\sum_{n\in H}u_n$, co wynika
bezpośrednio z definicji sum skończonych i jednoznaczności
słabej granicy. Obie wartości leżą w tej samej kuli,
stąd
$\|\sum_{n\in H}u_n\|\le2\varepsilon/3<\varepsilon$.
Otrzymaliśmy kryterium Cauchy'ego dla sum po dowolnych
skończonych podzbiorach ogona szeregu. W szczególności
dla każdego $M$ sumy $\sum_{n\in M,\ n\le N}u_n$
są ciągiem Cauchy'ego w normie, więc mają granicę w $X$.
Musi być ona równa $v_M$, ponieważ zbieżność w normie
implikuje słabą i słaba granica jest jednoznaczna.

Dla dowolnej permutacji indeksów sumy początkowe po
dostatecznie wielu wyrazach zawierają wszystkie indeksy
do $N_0$; pozostałe składniki mają normę sumy mniejszą
od $\varepsilon$. Suma szeregu w naturalnej kolejności
różni się od sumy do $N_0$ o wektor normy co najwyżej
$2\varepsilon/3$, przez przejście do granicy.
W istocie każda dostatecznie długa suma po permutacji
różni się od sumy do $N_0$ o wektor normy co najwyżej
$2\varepsilon/3$, na mocy tego samego oszacowania dla $H$.
Jej odległość od sumy szeregu jest zatem co najwyżej
$4\varepsilon/3$. Ponieważ $\varepsilon>0$ było dowolne,
każda permutacja daje tę samą granicę w normie.
To jest bezwarunkowa zbieżność.
\end{proof}

\begin{twierdzenie}[Miara generowana przez całkę Pettisa]
\label{thm:pettis-vector-measure15b}
Niech $(S,\Sigma,\mu)$ będzie zupełną przestrzenią z miarą
skończoną, $X$ przestrzenią Banacha nad $\K\in\{\R,\C\}$,
a $f:S\to X$ funkcją całkowalną w sensie Pettisa.
Wówczas $m_f:\Sigma\to X$, $m_f(A)=P_A(f)$,
jest przeliczalnie addytywną w normie miarą wektorową,
bezwzględnie ciągłą względem $\mu$.
Dla każdego ciągu parami rozłącznych $A_j\in\Sigma$ szereg
$\sum_jP_{A_j}(f)$ jest bezwarunkowo zbieżny w normie
i ma sumę $P_{\bigcup_j A_j}(f)$.
Wariacja $|m_f|(S)$ nie musi być skończona.
\end{twierdzenie}
\begin{proof}
Ustalmy dowolny ciąg parami rozłącznych zbiorów
$(A_j)_{j\ge1}$ z $\Sigma$ i niech $u_j=P_{A_j}(f)$.
Dla dowolnego $M\subseteq\N$
zbiór $B_M=\bigcup_{j\in M}A_j$ jest mierzalny, a
$v_M=P_{B_M}(f)$ istnieje na mocy całkowalności Pettisa.
Dla każdego $x'\in X'$ addytywność całki skalarnej daje
\[
 \lim_{N\to\infty}\sum_{\substack{j\in M\\j\le N}}x'(u_j)
 =\int_{B_M}x'(f(s))\,\mathrm d\mu(s)=x'(v_M).
\]
Wszystkie podszeregi mają zatem słabe granice w $X$.
Z twierdzenia Orlicza--Pettisa otrzymujemy bezwarunkową
zbieżność w normie i dla $M=\N$ żądaną addytywność.
Całka po zbiorze miary zero jest równa zeru, ponieważ
wszystkie jej skalaryzacje są zerowe; dotyczy to również
zbioru pustego. Brak skończonej wariacji wykazaliśmy
w \hyperref[ex:pettis-infinite-variation15b]{przykładzie z $\ell^2$}.
\end{proof}

\begin{twierdzenie}[Półwariacja miary Pettisa i ciągłość całek]
\label{thm:pettis-semivariation15b}
Niech $(S,\Sigma,\mu)$ będzie zupełną przestrzenią z miarą
skończoną, $X$ przestrzenią Banacha nad $\K\in\{\R,\C\}$,
a $f:S\to X$ funkcją całkowalną w sensie Pettisa.
Niech $m_f(A)=P_A(f)$ dla $A\in\Sigma$. Definiujemy
\emph{półwariację} miary $m_f$ na $A$ jako
$\|m_f\|_{\mathrm{sv}}(A)=
\sup_{\|x'\|\le1}|x'\circ m_f|(A)$, gdzie
$|x'\circ m_f|$ oznacza wariację miary skalarnej
$B\mapsto x'(m_f(B))$.
Wówczas
\[
 \|m_f\|_{\mathrm{sv}}(A)=q_f(A)
 =\sup_{\|\varphi\|_\infty\le1}\|P_A(\varphi f)\|,
\]
gdzie supremum jest brane po $\varphi\in L^\infty(S,\mu;\K)$,
a $q_f$ określono w
\hyperref[thm:vector-integral-multipliers15b]{twierdzeniu o mnożnikach}.
Ponadto
\[
 \lim_{\delta\downarrow0}
 \sup_{\substack{A\in\Sigma\\\mu(A)\le\delta}}q_f(A)=0.
\]
W szczególności dla $A_n,A\in\Sigma$ warunek
$\mu(A_n\triangle A)\to0$ implikuje
$\|P_{A_n}(f)-P_A(f)\|\to0$.
\end{twierdzenie}
\begin{proof}
Miara skalarna $x'\circ m_f$ ma gęstość $x'\circ f$
względem $\mu$. Jej wariacja na $A$ wynosi
$\int_A|x'(f)|\,\mathrm d\mu$: górne oszacowanie wynika
z nierówności trójkąta dla całki skalarnej, dolne
otrzymujemy przez aproksymację funkcji fazowej
$\overline{x'(f)}/|x'(f)|$ na zbiorze, na którym mianownik
jest niezerowy (przyjmujemy zero w pozostałych punktach),
prostymi funkcjami o module nie większym niż $1$.
Dla takiej funkcji prostej moduł całki iloczynu jest
nie większy od sumy modułów całek po składnikach podziału,
czyli od wariacji. Jednostajna aproksymacja fazy daje
w granicy całkę modułu. Biorąc supremum po $x'$,
otrzymujemy pierwszą równość.
Oszacowanie z twierdzenia o mnożnikach daje
$\sup_{\|\varphi\|_\infty\le1}\|P_A(\varphi f)\|\le q_f(A)$.
Dla ustalonego $x'$ z kuli jednostkowej wybieramy
$\varphi$ równą opisanej wyżej fazie; wtedy
$x'(P_A(\varphi f))=\int_A|x'(f)|\,\mathrm d\mu$.
Supremum po $x'$ daje nierówność przeciwną.

Dla $A\in\Sigma$ połóżmy
$p(A)=\sup\{\|m_f(B)\|:B\in\Sigma,\ B\subseteq A\}$.
Mamy $p(A)\le q_f(A)\le4p(A)$.
Pierwsza nierówność wynika z oszacowania całki.
Dla drugiej ustalmy $x'$ o normie co najwyżej $1$ i
oznaczmy $u=x'\circ f$. W przypadku zespolonym
$\int_A|u|\le\int_A|\operatorname{Re}u|
+\int_A|\operatorname{Im}u|$.
Całkę modułu części rzeczywistej rozbijamy na zbiory,
na których jest ona nieujemna i ujemna. Każda z dwóch
otrzymanych całek jest modułem części rzeczywistej
$x'(m_f(B))$ dla pewnego mierzalnego $B\subseteq A$,
zatem jest nie większa niż $p(A)$. Tak samo postępujemy
z częścią urojoną. W przypadku rzeczywistym wystarczają
dwa składniki. Biorąc supremum po $x'$, uzyskujemy
żądane oszacowanie.

Pokażemy, że jeżeli $E_j\downarrow\emptyset$, to
$p(E_j)\to0$. Z przeliczalnej addytywności $m_f$ w normie
wynika najpierw $m_f(C_j)\to0$ dla każdego
$C_j\downarrow\emptyset$: rozkład $C_1$ na parami
rozłączne różnice $C_j\setminus C_{j+1}$ daje szereg
zbieżny w normie, którego ogon od $j$ jest równy $m_f(C_j)$.
Gdyby $p(E_j)$ nie dążyło do zera, to z monotoniczności
$p$ istniałoby $\varepsilon>0$ takie, że dla każdego $j$
można wybrać $B\subseteq E_j$ mierzalny z $\|m_f(B)\|>2\varepsilon$.
Rekurencyjnie wybieramy indeksy $j_1<j_2<\cdots$ i
$B_k\subseteq E_{j_k}$ z $\|m_f(B_k)\|>2\varepsilon$.
Po wyborze $B_k$ dobieramy $j_{k+1}$ tak, aby
$\|m_f(B_k\cap E_{j_{k+1}})\|<\varepsilon$, co jest
możliwe dzięki ciągłości dla malejących zbiorów.
Zbiory $D_k=B_k\setminus E_{j_{k+1}}$ są parami rozłączne,
a ze skończonej addytywności i nierówności trójkąta
wynika $\|m_f(D_k)\|>\varepsilon$. Przeczy to
przeliczalnej addytywności w normie, bo wyrazy szeregu
$\sum_km_f(D_k)$ zbieżnego w normie muszą dążyć do zera.
Zatem $p(E_j)\to0$ i również $q_f(E_j)\to0$.

Jeżeli $E_j\downarrow E$ oraz $\mu(E)=0$, to
$q_f(E_j)=q_f(E_j\setminus E)$, ponieważ całki skalarne
po $E$ znikają. Stosujemy wynik do $E_j\setminus E$.
Gdyby żądana jednostajna ciągłość względem $\mu$ nie
zachodziła, istniałoby $\varepsilon>0$ i zbiory $A_j$
takie, że $\mu(A_j)\le2^{-j}$, lecz $q_f(A_j)\ge\varepsilon$.
Niech $E_j=\bigcup_{k\ge j}A_k$. Wtedy $E_j$ maleją,
$\mu(E_j)\le\sum_{k\ge j}2^{-k}\to0$, więc ich
przecięcie ma miarę zero. Zatem $q_f(E_j)\to0$,
wbrew $q_f(E_j)\ge q_f(A_j)\ge\varepsilon$.
Na koniec skończona addytywność daje
$P_{A_n}(f)-P_A(f)=P_{A_n\setminus A}(f)-P_{A\setminus A_n}(f)$.
Norma tej różnicy jest nie większa niż
$2q_f(A_n\triangle A)$, co kończy dowód.
\end{proof}

\begin{uwaga}[Półwariacja w przykładzie atomowym]
Dla miary $m(A)=(\chi_A(n)/n)_{n\ge1}$ z
\hyperref[ex:pettis-infinite-variation15b]{przykładu nieskończonej wariacji}
nierówność Cauchy'ego--Schwarza daje
\[
 \|m\|_{\mathrm{sv}}(A)
 =\sup_{\|a\|_{\ell^2}\le1}\sum_{n\in A}\frac{|a_n|}{n}
 =\left(\sum_{n\in A}\frac1{n^2}\right)^{1/2}.
\]
Równość zachodzi, gdy $a$ jest wektorem
$(\chi_A(n)/n)_{n\ge1}$ podzielonym przez jego normę
(dla $A=\emptyset$ obie strony są zerowe). W szczególności
półwariacja na $S$ jest skończona, choć całkowita wariacja
na $S$ jest nieskończona.
\end{uwaga}

\begin{twierdzenie}[Zbieżność skalarna i zbieżność w normie całek]
\label{thm:vector-integral-convergence15b}
Niech $(S,\Sigma,\mu)$ będzie zupełną przestrzenią z miarą
skończoną, $X,Y$ przestrzeniami Banacha nad
$\K\in\{\R,\C\}$, $(f_n)$ ciągiem funkcji całkowalnych
w sensie Dunforda z $S$ do $X$, a $(h_n)$ ciągiem funkcji
całkowalnych w sensie Gelfanda z $S$ do $Y'$.
Niech $f:S\to X$ i $h:S\to Y'$ będą funkcjami.
Jeżeli dla każdego $x'\in X'$ funkcje $x'\circ f_n$
zbiegają prawie wszędzie do $x'\circ f$ i istnieje
$g_{x'}\in L^1(S,\mu)$, $g_{x'}\ge0$, taka że
$|x'\circ f_n|\le g_{x'}$ prawie wszędzie dla każdego $n$,
to $f$ jest całkowalna w sensie Dunforda i dla każdego
$A\in\Sigma$ zachodzi
$D_A(f_n)\to D_A(f)$ w topologii $\sigma(X'',X')$.
Analogicznie, jeśli dla każdego $y\in Y$ funkcje
$h_n(\cdot)(y)$ zbiegają prawie wszędzie do $h(\cdot)(y)$
i są zmajoryzowane przez pewną nieujemną funkcję
$g_y\in L^1(S,\mu)$, to $h$ jest całkowalna w sensie
Gelfanda i $G_A(h_n)\to G_A(h)$ w $\sigma(Y',Y)$.
Jeżeli w pierwszym przypadku wszystkie $f_n$ oraz $f$
są całkowalne w sensie Pettisa, to zachodzi również
$P_A(f_n)\to P_A(f)$ w $\sigma(X,X')$.

Jeżeli ponadto $q_{f_n-f}(S)\to0$, to całki Dunforda
zbiegają w normie jednostajnie względem $A\in\Sigma$;
w przypadku Pettisa ten sam warunek daje zbieżność
w normie całek Pettisa. Dla Gelfanda analogicznym
warunkiem jest $r_{h_n-h}(S)\to0$.
Wystarczający warunek dla takiej zbieżności to
$\|f_n-f\|\to0$ prawie wszędzie i mierzalność tych norm
oraz oszacowanie $\|f_n-f\|\le g$ prawie wszędzie dla
wspólnej nieujemnej $g\in L^1(S,\mu)$; analogicznie dla
$h_n-h$.
\end{twierdzenie}
\begin{proof}
Dla ustalonego $x'$ zupełność miary zapewnia mierzalność
granicy prawie wszędzie $x'\circ f$, a skalarne twierdzenie
o zbieżności zmajoryzowanej daje jej całkowalność i
$\int_Ax'(f_n)\,\mathrm d\mu\to\int_Ax'(f)\,\mathrm d\mu$.
Z definicji całki Dunforda oznacza to zadeklarowaną
zbieżność słabą-*. Jeżeli całki Pettisa istnieją,
ta sama równość jest dokładnie zbieżnością słabą
ich wartości. Dla Gelfanda stosujemy ten argument
do każdego $y\in Y$. Zbiory wyjątkowe mogą zależeć
od testującego funkcjonału lub wektora, ponieważ
argument prowadzimy oddzielnie dla każdej skalaryzacji.

Wszystkie różnice należą do odpowiednich klas
całkowalności na mocy liniowości. Dla każdego $A$
z oszacowań całek i monotoniczności $q$ oraz $r$ mamy
$\|D_A(f_n)-D_A(f)\|\le q_{f_n-f}(S)$ oraz
$\|G_A(h_n)-G_A(h)\|\le r_{h_n-h}(S)$.
Dla Pettisa zachodzi identyczne oszacowanie z $q$,
ponieważ normę w $X$ wyznaczają funkcjonały z kuli
jednostkowej $X'$. Prawa strona nie zależy od $A$.
Przy dodatkowych założeniach o mierzalności i dominacji
normy różnicy zachodzi
$q_{f_n-f}(S)\le\int_S\|f_n-f\|\,\mathrm d\mu\to0$
z twierdzenia o zbieżności zmajoryzowanej; analogicznie
dla $r_{h_n-h}$.
\end{proof}

\begin{uwaga}[Zbieżność skalarna nie zapewnia zbieżności w normie]
Na jednopunktowej przestrzeni z miarą masy $1$ niech
$f_n(s)=e_n\in\ell^2$. Wszystkie funkcje są całkowalne
nawet w sensie Bochnera; dla każdego funkcjonału
$x'\in(\ell^2)'$ mamy $x'(e_n)\to0$ z reprezentacji
Riesza, a $|x'(e_n)|\le\|x'\|$. Ich całki Pettisa są
równe $e_n$ i zbiegają słabo do zera, lecz
$\|e_n\|=1$. Całki Dunforda są równe $J_{\ell^2}(e_n)$,
więc również nie zbiegają w normie do zera.
Dla Gelfanda przyjmujemy na tej samej przestrzeni
$h_n(s)=e_n\in(\ell^1)'=\ell^\infty$. Dla $y\in\ell^1$
mamy $h_n(s)(y)=y_n\to0$ i $|y_n|\le\|y\|_{\ell^1}$,
więc całki Gelfanda zbiegają słabo-* do zera, ale
ich normy są stale równe $1$. Pokazuje to konieczność
odróżnienia dominacji poszczególnych skalaryzacji od
warunku jednostajnego względem kuli jednostkowej.
\end{uwaga}

\begin{stwierdzenie}[Zależność całki Gelfanda od całki Dunforda i Pettisa]
\label{prop:gelfand-dunford-restriction15b}
Niech $(S,\Sigma,\mu)$ będzie zupełną przestrzenią
z miarą skończoną, $Y$ przestrzenią Banacha nad
$\K\in\{\R,\C\}$, a $h:S\to Y'$ funkcją całkowalną
w sensie Gelfanda. Jeśli dodatkowo $h$, traktowana jako
funkcja o wartościach w przestrzeni Banacha $Y'$, jest
całkowalna w sensie Dunforda, to dla każdego $A\in\Sigma$
\[
 G_A(h)=J_Y' D_A(h),
 \qquad J_Y':Y'''\to Y',\quad
 (J_Y'F)(y)=F(J_Yy).
\]
Jeżeli $h$ jest całkowalna w sensie Pettisa w $Y'$, to
$G_A(h)=P_A(h)$ dla wszystkich $A\in\Sigma$.
Odwrotne implikacje dotyczące całkowalności Pettisa
nie zachodzą w ogólności.
\end{stwierdzenie}
\begin{proof}
Zanurzenie kanoniczne $J_Y:Y\to Y''$ spełnia
$(J_Yy)(y')=y'(y)$. Dla każdego $y\in Y$ mamy
$D_A(h)(J_Yy)=\int_A(J_Yy)(h(s))\,\mathrm d\mu(s)
=\int_Ah(s)(y)\,\mathrm d\mu(s)=G_A(h)(y)$.
To daje pierwszą równość. Jeżeli istnieje całka Pettisa
$P_A(h)\in Y'$, to zastosowanie definicji do funkcjonału
$J_Yy\in Y''$ daje $P_A(h)(y)=G_A(h)(y)$
dla każdego $y$, czyli drugą równość. Kontrprzykład
poniżej pokazuje, że istnienie całki Gelfanda (nawet
wraz z całkowalnością w sensie Dunforda) nie zapewnia
całkowalności Pettisa.
\end{proof}

\begin{przyklad}[Całka Gelfanda niebędąca całką Pettisa]
\label{ex:gelfand-not-pettis15b}
Niech $S=\N=\{1,2,\ldots\}$, $\Sigma=\mathcal P(\N)$,
$\mu(\{n\})=2^{-n}$, $Y=\ell^1(\N;\K)$ i
$Y'=\ell^\infty(\N;\K)$, $\K\in\{\R,\C\}$.
W \hyperref[ex:gelfand-weakstar15b]{przykładzie słabej-* addytywności}
rozważaliśmy $h(n)=2^ne_n$ i otrzymaliśmy $G_A(h)=\chi_A$.
Funkcja $h$ jest nawet całkowalna w sensie Dunforda
w przestrzeni $\ell^\infty$. Istotnie, jest silnie
mierzalna przez skończone obcięcia, a dla każdego
$\Phi\in(\ell^\infty)'$ oraz $N\ge1$ wybieramy skalary
$c_n$ o module co najwyżej $1$, dla których
$c_n\Phi(e_n)=|\Phi(e_n)|$. Wówczas
$\|\sum_{n=1}^Nc_ne_n\|_\infty\le1$, stąd
$\sum_{n=1}^N|\Phi(e_n)|
=|\Phi(\sum_{n=1}^Nc_ne_n)|\le\|\Phi\|$. Przejście do granicy
daje $\int_S|\Phi(h(n))|\,\mathrm d\mu(n)
=\sum_n|\Phi(e_n)|<\infty$.

Gdyby $h$ była całkowalna w sensie Pettisa, to jej
całka po $S$ musiałaby być równa $G_S(h)=(1,1,\ldots)$
na mocy poprzedniego stwierdzenia. Niech $c\subset\ell^\infty$
oznacza podprzestrzeń ciągów zbieżnych. Funkcjonał
$L:c\to\K$, $L((b_n))=\lim_n b_n$, jest liniowy
i ograniczony normą $1$. Twierdzenie Hahna--Banacha
daje jego ciągłe przedłużenie $\Phi$ na $\ell^\infty$.
Mamy $\Phi(e_n)=0$ dla każdego $n$, ponieważ
$e_n\in c$ ma granicę zero, natomiast
$\Phi((1,1,\ldots))=1$. Zatem
$\int_S\Phi(h(n))\,\mathrm d\mu(n)=0\ne1
=\Phi(G_S(h))$, co przeczy definicji całki Pettisa.
\end{przyklad}

\subsection{Twierdzenia Fubiniego dla całek słabych}

W twierdzeniu Fubiniego dla całki Bochnera silna mierzalność
pozwala wybrać jeden zbiór wyjątkowy dla mierzalności
przekrojów, a całkowalność normy na produkcie zapewnia
całkowalność normy prawie wszystkich przekrojów dzięki
skalarnemu twierdzeniu Tonellego.
Przy całkach słabych zbiory wyjątkowe ze skalarnego
twierdzenia Fubiniego mogą zależeć od funkcjonału.
Nie wolno bez dodatkowego argumentu zastąpić ich
nieprzeliczalnej rodziny jednym zbiorem miary zero.

\begin{twierdzenie}[Twierdzenie Fubiniego dla całki Gelfanda]
\label{thm:gelfand-fubini15b}
Niech $(S,\Sigma,\mu)$ i $(T,\mathcal T,\nu)$ będą
zupełnymi przestrzeniami z miarami skończonymi,
$\rho$ uzupełnieniem miary $\mu\otimes\nu$, a $Y$
przestrzenią Banacha nad $\K\in\{\R,\C\}$.
Niech $h:S\times T\to Y'$ będzie całkowalna w sensie
Gelfanda względem $\rho$. Załóżmy, że istnieją
mierzalne zbiory zerowe $N_S\subseteq S$, $N_T\subseteq T$
takie, że dla każdego $s\notin N_S$ funkcja $h(s,\cdot)$
jest całkowalna w sensie Gelfanda względem $\nu$,
a dla każdego $t\notin N_T$ funkcja $h(\cdot,t)$ jest
całkowalna w sensie Gelfanda względem $\mu$.
Dla $B\in\mathcal T$ i $A\in\Sigma$ określmy
\[
 H_B(s)=\int_B^{\mathrm G}h(s,t)\,\mathrm d\nu(t),
 \qquad K_A(t)=\int_A^{\mathrm G}h(s,t)\,\mathrm d\mu(s),
\]
przyjmując wartości zero odpowiednio na $N_S$ i $N_T$.
Wówczas $H_B:S\to Y'$ i $K_A:T\to Y'$ są całkowalne
w sensie Gelfanda oraz
\[
 \int_{A\times B}^{\mathrm G}h\,\mathrm d\rho
 =\int_A^{\mathrm G}H_B\,\mathrm d\mu
 =\int_B^{\mathrm G}K_A\,\mathrm d\nu .
\]
\end{twierdzenie}
\begin{proof}
Ustalmy $y\in Y$ oraz $B\in\mathcal T$.
Funkcja skalarna $u_y(s,t)=h(s,t)(y)$ jest całkowalna
względem $\rho$. Skalarne twierdzenie Fubiniego
dla uzupełnionej miary produktowej daje całkowalność
przekroju i mierzalność funkcji jego całek dla prawie
każdego $s$. Można stosować to twierdzenie, ponieważ
$u_y$ ma produktowo mierzalnego przedstawiciela
równego jej prawie wszędzie. Zbiór, na którym funkcja
i przedstawiciel się różnią, zawiera się w produktowo
mierzalnym zbiorze miary zero. Twierdzenie Tonellego
dla jego funkcji charakterystycznej daje przekroje
miary zero dla prawie każdego parametru; zupełność miar
czynników pozwala wówczas zastąpić przekroje przedstawiciela
przekrojami pierwotnej funkcji.
Na $S\setminus N_S$ z definicji całki Gelfanda zachodzi
$H_B(s)(y)=\int_B u_y(s,t)\,\mathrm d\nu(t)$.
Funkcja $H_B(\cdot)(y)$ jest więc prawie wszędzie
równa mierzalnej funkcji skalarnej z twierdzenia Fubiniego.
Zupełność $\mu$ daje jej mierzalność, a ponadto
\[
 \int_S|H_B(s)(y)|\,\mathrm d\mu(s)
 \le\int_{S\times B}|u_y|\,\mathrm d\rho<\infty.
\]
Ponieważ $y$ było dowolne, $H_B$ jest całkowalna
w sensie Gelfanda. Dla każdego $A\in\Sigma$ mamy
\[
 \left(\int_A^{\mathrm G}H_B\,\mathrm d\mu\right)(y)
 =\int_A\int_B u_y(s,t)\,\mathrm d\nu(t)\,\mathrm d\mu(s)
 =\int_{A\times B}u_y\,\mathrm d\rho.
\]
Prawa strona jest wartością całki Gelfanda funkcji $h$
po $A\times B$ na wektorze $y$. Równość dla wszystkich
$y\in Y$ daje równość funkcjonałów. Zamiana ról
przestrzeni $S$ i $T$ daje drugi porządek całkowania.
\end{proof}

\begin{twierdzenie}[Twierdzenie Fubiniego dla całki Dunforda]
\label{thm:dunford-fubini15b}
Niech $(S,\Sigma,\mu)$ i $(T,\mathcal T,\nu)$ będą
zupełnymi przestrzeniami z miarami skończonymi,
$\rho$ uzupełnieniem $\mu\otimes\nu$, a $X$ przestrzenią
Banacha nad $\K\in\{\R,\C\}$.
Niech $f:S\times T\to X$ będzie całkowalna w sensie
Dunforda względem $\rho$. Załóżmy, że poza mierzalnymi
zbiorami miary zero parametrów jej przekroje
$f(s,\cdot)$ i $f(\cdot,t)$ są całkowalne w sensie
Dunforda względem odpowiednich miar.
Dla $A\in\Sigma$ i $B\in\mathcal T$ połóżmy
$U_B(s)=D_B(f(s,\cdot))\in X''$ oraz
$V_A(t)=D_A(f(\cdot,t))\in X''$, z wartością zero
na odpowiednich zbiorach wyjątkowych.
Wówczas $U_B$ i $V_A$ są całkowalne w sensie Gelfanda
w przestrzeni sprzężonej $(X')'=X''$ oraz
\[
 D_{A\times B}(f)=
 \int_A^{\mathrm G}U_B\,\mathrm d\mu
 =\int_B^{\mathrm G}V_A\,\mathrm d\nu.
\]
Nie stwierdzamy tu całkowalności Dunforda funkcji
$U_B$ i $V_A$ względem wszystkich funkcjonałów z $X'''$.
\end{twierdzenie}
\begin{proof}
Zanurzenie kanoniczne $J_X:X\to X''$ daje funkcję
$h=J_X\circ f:S\times T\to (X')'$. Dla każdego $x'\in X'$
mamy $h(s,t)(x')=x'(f(s,t))$. Mierzalność i całkowalność
tych funkcji skalarnych są dokładnie założeniami
o całkowalności Dunforda funkcji $f$. Stąd $h$ jest
całkowalna w sensie Gelfanda z przestrzenią testującą
$Y=X'$, a dla dowolnego mierzalnego $C$ w uzupełnieniu
produktu jej całka Gelfanda po $C$ jest równa $D_C(f)$.
Ten sam argument dla przekrojów identyfikuje ich całki
Gelfanda z $U_B$ i $V_A$. Twierdzenie Fubiniego
dla Gelfanda zastosowane do $h$ daje wszystkie
stwierdzone wnioski.
\end{proof}

\begin{twierdzenie}[Twierdzenie Fubiniego dla całki Pettisa]
\label{thm:pettis-fubini15b}
Niech $(S,\Sigma,\mu)$ i $(T,\mathcal T,\nu)$ będą
zupełnymi przestrzeniami z miarami skończonymi,
$\rho$ uzupełnieniem $\mu\otimes\nu$, a $X$ przestrzenią
Banacha nad $\K\in\{\R,\C\}$.
Niech $f:S\times T\to X$ będzie całkowalna w sensie
Pettisa względem $\rho$. Załóżmy ponadto, że poza
mierzalnymi zbiorami miary zero parametrów przekroje
$f(s,\cdot)$ i $f(\cdot,t)$ są całkowalne w sensie
Pettisa względem odpowiednich miar.
Dla $A\in\Sigma$ i $B\in\mathcal T$ definiujemy
\[
 F_B(s)=\int_B^{\mathrm P}f(s,t)\,\mathrm d\nu(t),
 \qquad L_A(t)=\int_A^{\mathrm P}f(s,t)\,\mathrm d\mu(s),
\]
z wartością zero na odpowiednich zbiorach wyjątkowych.
Wówczas $F_B$ i $L_A$ są całkowalne w sensie Pettisa
i zachodzi
\[
 \int_{A\times B}^{\mathrm P}f\,\mathrm d\rho
 =\int_A^{\mathrm P}F_B\,\mathrm d\mu
 =\int_B^{\mathrm P}L_A\,\mathrm d\nu.
\]
\end{twierdzenie}
\begin{proof}
Całkowalność Pettisa implikuje całkowalność Dunforda,
a dla przekrojów mamy $U_B=J_X\circ F_B$ i
$V_A=J_X\circ L_A$ w oznaczeniach poprzedniego twierdzenia.
Zatem z całkowalności Gelfanda funkcji $U_B$ wynika,
że dla każdego $x'\in X'$ funkcja
$x'(F_B)=U_B(\cdot)(x')$ jest mierzalna i całkowalna
względem $\mu$. Nie wystarcza to jeszcze do całkowalności
Pettisa: potrzebna jest reprezentacja całek po wszystkich
mierzalnych podzbiorach $S$.

Niech więc $C\in\Sigma$ będzie dowolny. Z twierdzenia
Fubiniego dla Dunforda i całkowalności Pettisa na
produkcie otrzymujemy
\[
 D_C(F_B)=D_{C\times B}(f)
 =J_X\left(\int_{C\times B}^{\mathrm P}f\,\mathrm d\rho\right).
\]
Pierwsza równość wynika z testowania obu stron
na każdym $x'\in X'$; druga z definicji Pettisa.
W szczególności $D_C(F_B)\in J_X(X)$ dla każdego $C$.
Kryterium reprezentowalności całki Dunforda daje
całkowalność Pettisa funkcji $F_B$ i żądaną równość
dla $C=A$. Dla $L_A$ postępujemy analogicznie,
wykorzystując wszystkie prostokąty $A\times D$,
$D\in\mathcal T$.
\end{proof}

\begin{stwierdzenie}[Warunki zapewniające całkowalność przekrojów]
\label{prop:weak-fubini-separable15b}
Niech $(S,\Sigma,\mu)$, $(T,\mathcal T,\nu)$ będą
zupełnymi przestrzeniami z miarami skończonymi,
$\rho$ uzupełnieniem $\mu\otimes\nu$, a $Y$ ośrodkową
przestrzenią Banacha nad $\K\in\{\R,\C\}$.
Niech $h:S\times T\to Y'$ będzie słabo-* mierzalna
względem $\rho$. Załóżmy, że istnieje nieujemna
$\rho$-mierzalna funkcja $g\in L^1(\rho)$ taka, że
$\|h(s,t)\|\le g(s,t)$ dla $\rho$-prawie wszystkich
$(s,t)$. Wtedy $h$ i prawie wszystkie przekroje
w obu kierunkach są całkowalne w sensie Gelfanda,
więc zachodzi twierdzenie Fubiniego dla Gelfanda.
Nie trzeba osobno zakładać mierzalności normy $h$:
wynika ona z ośrodkowości $Y$ i słabej-* mierzalności $h$.

Jeżeli $X$ jest przestrzenią Banacha nad $\K$
o ośrodkowej przestrzeni sprzężonej $X'$,
$f:S\times T\to X$ jest słabo mierzalna względem $\rho$
i $\|f\|\le g$ prawie wszędzie dla takiego $g$,
to zachodzi analogiczny wniosek dla Dunforda.
Jeżeli dodatkowo $X$ jest refleksywna, to zachodzi
również wniosek dla Pettisa.
\end{stwierdzenie}
\begin{proof}
Dla każdego $y\in Y$ nierówność
$|h(s,t)(y)|\le\|y\|g(s,t)$ prawie wszędzie zapewnia
całkowalność skalarną na produkcie. Wybierzmy gęsty
ciąg $(y_k)$ w domkniętej kuli jednostkowej $Y$.
Dla każdego $(s,t)$ ciągłość funkcjonału $h(s,t)$ daje
$\|h(s,t)\|=\sup_k|h(s,t)(y_k)|$, więc norma jest
mierzalna jako supremum przeliczalnej rodziny funkcji
mierzalnych. Dla każdego $k$ skalarna
mierzalność na produkcie zapewnia mierzalność
przekrojów $h(s,\cdot)(y_k)$ dla prawie każdego $s$
i $h(\cdot,t)(y_k)$ dla prawie każdego $t$.
Można to uzasadnić przez wybór produktowo mierzalnego
przedstawiciela i zastosowanie skalarnego Fubiniego
do funkcji charakterystycznej mierzalnego zbioru
zerowego, zawierającego zbiór, na którym przedstawiciele
się różnią.

Ponieważ rodzina $(y_k)$ jest przeliczalna, można
usunąć jeden zbiór miary zero w $S$, poza którym
wszystkie te przekroje są mierzalne. Powiększamy
ten zbiór o zbiór wyjątkowy ze skalarnego Fubiniego
dla $g$ oraz o zbiór parametrów, dla których przekrój
zbioru, na którym nie zachodzi dominacja, nie jest
zerowy. Ściślej, ten ostatni zbiór w produkcie
zawieramy w produktowo mierzalnym zbiorze miary zero
i stosujemy do niego skalarne twierdzenie Fubiniego.
W konsekwencji dla każdego pozostałego $s$ funkcja
$g(s,\cdot)$ jest mierzalna i całkowalna, a
$\|h(s,t)\|\le g(s,t)$ dla $\nu$-prawie każdego $t$.

Ustalmy taki $s$ i dowolny $y\in Y$. Dla $y=0$
skalaryzacja jest zerowa. Dla $y\ne0$ wybieramy
dla każdego $j\ge1$ indeks $k_j$ taki, że
$\|y_{k_j}-y/\|y\|\|<1/j$.
Dla każdego $t\in T$ ciągłość funkcjonału $h(s,t)$
daje $\|y\|h(s,t)(y_{k_j})\to h(s,t)(y)$.
Funkcja ta jest zatem mierzalna jako granica funkcji
mierzalnych. Jej moduł jest prawie wszędzie ograniczony
przez $\|y\|g(s,\cdot)$, więc jest całkowalna.
To dowodzi całkowalności Gelfanda przekroju dla każdego
$s$ poza jednym zbiorem zerowym, niezależnym od $y$.
Zamiana ról $S$ i $T$ daje drugi kierunek.

Dla Dunforda stosujemy powyższy argument do
$Y=X'$ i $h=J_X\circ f$, zauważając, że
$\|J_X(f(s,t))\|=\|f(s,t)\|$. Jeżeli $X$ jest
refleksywna, to $J_X(X)=X''$, zatem każda całka
Dunforda na produkcie i w przekrojach jest reprezentowana
przez wektor z $X$. Wtedy stosujemy twierdzenie
Fubiniego dla Pettisa.
\end{proof}

\begin{przyklad}[Całkowalność na produkcie nie gwarantuje całkowalności przekrojów]
\label{ex:weak-fubini-failure15b}
Niech $S=T=[0,1]$ z zupełną miarą Lebesgue'a
$\lambda$, $\rho$ uzupełnieniem $\lambda\otimes\lambda$,
a $X=\ell^2([0,1];\K)$, $\K\in\{\R,\C\}$.
Dla $s\in[0,1]$ oznaczmy przez $e_s$ wektor jednostkowy
o jedynej niezerowej współrzędnej na indeksie $s$.
Definiujemy $k(0)=0$, $k(t)=1/t$ dla $t>0$, oraz
$f(s,t)=k(t)e_s$.

Dla każdego $a\in X$ zbiór $C_a=\{s:a_s\ne0\}$
jest co najwyżej przeliczalny (jak wykazaliśmy
w \hyperref[ex:pettis-not-bochner15b]{przykładzie nieośrodkowej przestrzeni Hilberta}).
Przez twierdzenie Riesza wszystkie skalaryzacje $f$
mają postać $k(t)\overline{a_s}$ i znikają poza
$C_a\times T$. Ten zbiór jest produktowo mierzalny
i ma miarę $\rho$ równą zero. Zupełność $\rho$
daje mierzalność skalaryzacji, a ponieważ jest ona
równa zeru prawie wszędzie, jej całka modułu wynosi zero.
Dla dowolnego mierzalnego $C$ w uzupełnieniu produktu
wektor zero reprezentuje wszystkie całki skalarne
po $C$. Zatem $f$ jest całkowalna w sensie Pettisa,
a więc także Dunforda, na produkcie, i wszystkie
jej całki Pettisa są zerowe.

Ustalmy jednak dowolne $s\in S$. Funkcjonał
$x_s'(v)=\langle v,e_s\rangle$ spełnia
$x_s'(f(s,t))=k(t)$, a $\int_T k\,\mathrm d\lambda=\infty$.
Dlatego żaden przekrój $f(s,\cdot)$ nie jest całkowalny
nawet w sensie Dunforda. Zbiór złych parametrów
jest całym $S$, chociaż dla każdego ustalonego
funkcjonału jego zbiory wyjątkowe mogą być zerowe.

Dla ustalonego $t$ przekrój $f(\cdot,t)$ jest natomiast
całkowalny w sensie Pettisa z całką zero na każdym
mierzalnym podzbiorze $S$: wszystkie skalaryzacje
znikają poza przeliczalnymi zbiorami. Jeden porządek
całkowania jest więc możliwy, drugi nie istnieje.
Przyjmując $Y=X'$ oraz $h=J_X\circ f:S\times T\to Y'$,
otrzymujemy również kontrprzykład dla Gelfanda.
Istotnie, dla każdego $x'\in Y=X'$ zachodzi
$h(s,t)(x')=x'(f(s,t))$. Wszystkie skalaryzacje
funkcji $h$ na produkcie są więc całkowalne,
lecz dla każdego $s$ skalaryzacja przekroju przez
$x_s'\in Y$ nie jest całkowalna. Konstrukcja działa
nad obydwoma ciałami skalarów i nie wymaga identyfikacji
przestrzeni Hilberta z jej przestrzenią sprzężoną.
\end{przyklad}

\subsection{Własność Radona--Nikodyma}

\begin{definicja}[Własność Radona--Nikodyma]
\label{def:rnp15b}
Przestrzeń Banacha $X$ nad $\K\in\{\R,\C\}$ ma
\emph{własność Radona--Nikodyma}
(RNP), jeżeli dla każdej skończonej przestrzeni z miarą
$(S,\Sigma,\mu)$ oraz każdej miary wektorowej
$m:\Sigma\to X$ o skończonej całkowitej wariacji,
bezwzględnie ciągłej względem $\mu$ (tzn. $m(A)=0$ dla
każdego $A\in\Sigma$ spełniającego $\mu(A)=0$), istnieje funkcja
$f\in L^1(S,\mu;X)$ taka, że
$m(A)=\int_Af(s)\,\mathrm d\mu(s)$ dla wszystkich $A\in\Sigma$.
Przy mierze niezupełnej całkowalność w sensie Bochnera
rozumiemy względem jej uzupełnienia. Jak uzasadnimy poniżej,
każda taka funkcja ma przedstawiciela będącego punktową
granicą funkcji prostych mierzalnych względem $\Sigma$;
dla tego przedstawiciela całki norm można obliczać
względem pierwotnej miary $\mu$.
Można też ograniczyć definicję RNP do miar zupełnych:
miara wektorowa $m$ rozszerza się na uzupełnienie przez
$\bar m(A)=m(B)$, gdy $B\in\Sigma$ i różnica symetryczna
$A\triangle B$ jest zawarta w mierzalnym zbiorze miary zero.
Jest to niezależne od $B$: dla dwóch przedstawicieli
$B,C\in\Sigma$ ich różnica symetryczna ma miarę zero,
więc $m(B\setminus C)=m(C\setminus B)=0$.
Niech $(A_j)$ będzie ciągiem parami rozłącznych zbiorów
z uzupełnienia, a $B_j\in\Sigma$ ich przedstawicielami.
Zbiory $C_j=B_j\setminus\bigcup_{i<j}B_i$ są parami
rozłączne, a $C_j$ jest nadal przedstawicielem $A_j$,
gdyż wszystkie przecięcia $B_j\cap B_i$ dla $i\ne j$
mają miarę zero. Ponadto $\bigcup_j C_j$ jest przedstawicielem
$\bigcup_j A_j$, zatem przeliczalna addytywność $m$ daje
$\bar m(\bigcup_j A_j)=\sum_j m(C_j)=\sum_j\bar m(A_j)$
w normie. Dla skończonych podziałów $S$ taki sam wybór
rozłącznych przedstawicieli daje
$\sum_j\|\bar m(A_j)\|=\sum_j\|m(C_j)\|\le |m|(S)$.
Ponieważ $\bar m$ rozszerza $m$, zachodzi
$|\bar m|(S)=|m|(S)<\infty$. Rozszerzenie jest bezwzględnie
ciągłe względem uzupełnionej miary, bo przedstawiciel
zbioru miary zero również ma miarę zero.
Gęstość względem uzupełnienia ma przedstawiciela
silnie mierzalnego względem $\Sigma$: zastępujemy zbiory
w ciągu funkcji prostych mierzalnymi przedstawicielami,
a granicę określamy na mierzalnym zbiorze punktów,
w których ciąg jest Cauchy'ego, i jako zero poza nim.
Mnożąc nowe funkcje proste przez funkcję charakterystyczną
tego zbioru, uzyskujemy zbieżność punktową do przedstawiciela
na całym $S$. Przedstawiciel jest równy pierwotnej gęstości
prawie wszędzie i ma tę samą całkę normy. Przybliżenia
w normie $L^1$ można otrzymać przez przycięcie funkcji
prostych jak w dowodzie kryterium całkowalności Bochnera.
Jeżeli $X$ ma RNP dla miar zupełnych, zastosowanie tej
własności do $\bar m$ daje więc gęstość dla pierwotnej
miary $m$; implikacja odwrotna jest natychmiastowa.
\end{definicja}

\begin{przyklad}[Brak RNP w przestrzeni funkcji ciągłych]
\label{ex:rnp-counterexample15b}
Niech $X=C([0,1])$ z normą supremum i niech $\mu$ będzie
miarą Lebesgue'a na $[0,1]$. Definiujemy
$m(A)(t)=\mu(A\cap[0,t])$ dla $t\in[0,1]$.
Dla $0\le u\le t\le1$ mamy
$|m(A)(t)-m(A)(u)|=\mu(A\cap(u,t])\le t-u$,
więc $m(A)$ jest ciągła. Skończona addytywność $m$ i
równość $m(\emptyset)=0$ wynikają ze skalarnej addytywności
$\mu$. Dla parami rozłącznych $A_n$ zachodzi
$\|m(\bigcup_{n>N}A_n)\|_\infty
=\mu(\bigcup_{n>N}A_n)\to0$,
a skończona addytywność zastosowana do sumy wszystkich
$A_n$ rozbitej na pierwszych $N$ zbiorów i ogon daje
addytywność przeliczalną w normie. Zatem $m$ jest miarą
wektorową. Co więcej, $\|m(A)\|_\infty=\mu(A)$, więc
$|m|(A)=\mu(A)$; wariacja jest skończona i
bezwzględnie ciągła względem $\mu$.

Gdyby $m(A)=\int_Af(s)\,\mathrm d\mu(s)$ dla pewnej funkcji
$f\in L^1([0,1];C([0,1]))$, to dla każdego wymiernego
$q\in[0,1]$ i wszystkich zbiorów mierzalnych $A$
mielibyśmy
$\int_A f(s)(q)\,\mathrm d\mu(s)
=\int_A\chi_{[0,q]}(s)\,\mathrm d\mu(s)$.
Stąd $f(s)(q)=\chi_{[0,q]}(s)$ dla prawie wszystkich $s$.
Po usunięciu jednego wspólnego zbioru miary zero równość
zachodzi dla wszystkich wymiernych $q$ jednocześnie.
Dla takiego $s\in(0,1)$ wybieramy ciągi wymiernych
$q_n<s<r_n$ zbieżne do $s$. Wówczas $f(s)(q_n)=0$
i $f(s)(r_n)=1$ dla każdego $n$. Ciągłość $f(s)$
w punkcie $s$ wymusza jednocześnie $f(s)(s)=0$ i $f(s)(s)=1$.
Otrzymaliśmy sprzeczność: przestrzeń $C([0,1])$ nie ma RNP.
\end{przyklad}

\begin{uwaga}[Rozszerzenie na przestrzenie lokalnie wypukłe]
\label{rem:lcs-pettis15b}
Dla przestrzeni lokalnie wypukłej Hausdorffa $E$ sensowna
pozostaje \emph{całka Pettisa w $E$}: dla funkcji $f:S\to E$
wymagamy całkowalności wszystkich skalaryzacji $\ell\circ f$,
$\ell\in E'$, i dla każdego $A\in\Sigma$ istnienia $v_A\in E$
takiego, że $\ell(v_A)=\int_A\ell(f(s))\,\mathrm d\mu(s)$.
Rozdzielanie punktów przez $E'$ zapewnia jednoznaczność
$v_A$. Jeżeli $T:E\to F$ jest ciągłym odwzorowaniem liniowym
między przestrzeniami lokalnie wypukłymi Hausdorffa,
to $Tf$ jest całkowalna w sensie Pettisa w $F$ i
$T(v_A)=\int_A^{\mathrm P} T(f(s))\,\mathrm d\mu(s)$:
dla każdego $\varphi\in F'$ wystarczy zastosować
$\varphi\circ T\in E'$. Nie zakładamy jednak, że sama
zupełność $E$ wymusza istnienie całek Pettisa wszystkich
funkcji słabo mierzalnych: jest to własność mocniejsza.
\end{uwaga}

\begin{uwaga}[Zakres teorii]
W osobnym skrypcie teorii miary rozwijamy skalarne podstawy
całkowania, natomiast tu wyjaśniliśmy różnice pomiędzy
całkowaniem w normie, całkowaniem skalarnym z wartościami
w przestrzeni podwójnie sprzężonej i całkowaniem słabym-* do przestrzeni sprzężonej.
Równania różniczkowe wykorzystujące te konstrukcje należą
do osobnego skryptu PDE. Ujęcie całek dla ogólnych przestrzeni
liniowo-topologicznych nie musi być oparte na seminormach;
bez lokalnej wypukłości argumenty Hahna--Banacha nie są
automatycznie dostępne.
\end{uwaga}

\addtocounter{section}{-1}
\endgroup
